\documentclass[11pt,a4paper]{article}
\usepackage[T1]{fontenc}
\usepackage[utf8]{inputenc}
\usepackage{amsmath,amssymb,amsthm}
\newtheorem{proposition}{Proposition}[section]
\usepackage{newpxtext,newpxmath}
\usepackage{microtype}
\usepackage[margin=25mm,top=25mm,bottom=25mm]{geometry}
\usepackage[numbers,sort&compress]{natbib}
\usepackage{longtable,booktabs,array,calc}
\usepackage{fancyhdr,needspace}
\usepackage{xurl}
\usepackage[hidelinks]{hyperref}
\usepackage{bookmark}
\providecommand{\tightlist}{\setlength{\itemsep}{0pt}\setlength{\parskip}{0pt}}
\providecommand{\passthrough}[1]{#1}
\setlength{\emergencystretch}{2em}
\setlength{\parindent}{0pt}
\setlength{\parskip}{.45em}
\setlength{\headheight}{14pt}
\linespread{1.025}
\widowpenalty=10000
\clubpenalty=10000
\displaywidowpenalty=10000
\setcounter{tocdepth}{1}
\pagestyle{fancy}
\fancyhf{}
\fancyhead[L]{\small\itshape Seven Coincidences: proofs and computational checks}
\fancyfoot[C]{\thepage}
\renewcommand{\headrulewidth}{0pt}
\hypersetup{pdftitle={Proofs and Computational Checks for Seven Coincidences and Sixty Symmetries}}
\begin{document}
\thispagestyle{plain}
\begin{center}
{\Large\bfseries Proofs and Computational Checks}\\[6pt]
{\large for Seven Coincidences and Sixty Symmetries}\\[6pt]
Ernesto Lupercio and Evgeniya Shaydurova\\[3pt]
{\small Supplement to \emph{Seven Coincidences and Sixty Symmetries}}
\end{center}
\medskip
We prove the seven $A_5$-identities, enumerate their realizations, and develop the spinorial, arithmetic, Klein, modular, and Conway constructions used in the article. The accompanying programs carry out the finite character, partition and matrix calculations in exact arithmetic.

Unless otherwise stated, representations are finite-dimensional and complex. The symbols \(\mathbf1,A,B,C,E,P,Q,S\) have the same meaning as in the article. We write \(U\cong V\) when there is an equivariant linear isomorphism \(U\to V\). The construction of the Klein intertwiner requires choices of invariant splittings and identifications of multiplicity spaces.

Baez's question, and the spinorial explanation attributed there to Paul Schwahn, provide the immediate starting point \citep{Baez2026Coincidences}. We study representations realizing these seven equalities.

\tableofcontents

\clearpage

\appendix
\section{Finite representations and exterior-power identities}\label{sec:finite}

\subsection{Modules and the character criterion}\label{subsec:A1}

Let \(G=A_5\), acting on five letters. The module \(A\) is the zero-sum part of the five-letter permutation representation. Let \(H\) be a subgroup generated by a five-cycle and \(N=N_G(H)\) its order-ten normalizer. The coset spaces \(G/H\) and \(G/N\) have twelve and six elements and give \(Q\) and \(P\), respectively. The map \(G/H\to G/N\) has two-element fibers, and the involution interchanging the two points of each fiber splits the function space into its even and odd subspaces, giving \(Q=P\oplus C\). Let \(B\) be the zero-sum invariant complement to the constant functions in \(P\). These constructions give invariant direct-sum decompositions of the permutation modules. The edge module \(E\) is the permutation module on two-element subsets of the five letters, and \(S=2B\oplus C\).

For a finite-order transformation with eigenvalues \(\lambda_i\),

\[
\sum_{r\geq0}\chi_{\Lambda^rV}(g)t^r=\prod_i(1+t\lambda_i).
\]

Logarithmic differentiation gives Newton's recurrence

\[
r e_r=\sum_{i=1}^r(-1)^{i-1}\chi_V(g^i)e_{r-i},\qquad e_0=1.
\]

A permutation cycle of length \(\ell\) contributes \(1-(-t)^\ell\). The sixty even permutations, together with the coset and edge actions above, determine every character in the table. We evaluate each row in two equivalent ways, by the cycle-product formula and by Newton's recurrence. For an invariant decomposition \(V=U\oplus W\), the determinant polynomials satisfy \(\det(1+tg|_V)=\det(1+tg|_U)\det(1+tg|_W)\), so either factor can be recovered by division.

\paragraph{Why equality of characters is sufficient.}
Let \(G\) be any finite group and let \(W\) be a finite-dimensional complex representation. Starting from a positive-definite Hermitian inner product \(\langle\ ,\ \rangle_0\), define
\[
\langle u,v\rangle_G=\frac1{|G|}\sum_{g\in G}
       \langle\rho_W(g)u,\rho_W(g)v\rangle_0.
\]
This is positive definite and invariant, since multiplication by any fixed group element permutes the summation indices. If \(U\subset W\) is invariant, its orthogonal complement is invariant as well: for \(v\in U^\perp\), \(u\in U\), one has \(\langle gv,u\rangle_G=\langle v,g^{-1}u\rangle_G=0\). Splitting off a proper invariant subspace and proceeding by induction on dimension decomposes \(W\) into irreducible summands, where irreducible means having no invariant subspaces other than zero and the whole space.

If \(T,R\) are irreducible, an equivariant map \(T\to R\) has invariant kernel and image, and is consequently either zero or an isomorphism. An equivariant endomorphism of \(T\) is scalar: choose a complex eigenvalue \(\lambda\); the equivariant map \(A-\lambda I\) is not invertible, and hence is zero. It follows that the dimension of the space \(\operatorname{Hom}_G(T,W)\) is the multiplicity of \(T\) in an irreducible decomposition of \(W\).

That dimension can be computed from the characters. On the vector space of all linear maps \(A:T\to W\), average the group action:
\[
\mathcal P(A)=\frac1{|G|}\sum_{g\in G}
       \rho_W(g)A\rho_T(g)^{-1}.
\]
Changing the summation index shows that \(\mathcal P(A)\) is equivariant for every linear map \(A:T\to W\), and averaging fixes every equivariant map; thus \(\mathcal P\) is a projection onto \(\operatorname{Hom}_G(T,W)\). Its trace is therefore the dimension of its image. For matrices \(B,C\), the operator \(A\mapsto BAC\) on this space has trace \((\operatorname{Tr}B)(\operatorname{Tr}C)\), as is seen on the matrix-unit basis. Thus
\[
\dim\operatorname{Hom}_G(T,W)
=\frac1{|G|}\sum_{g\in G}\chi_W(g)\chi_T(g^{-1})
=\frac1{|G|}\sum_{g\in G}\chi_W(g)\overline{\chi_T(g)}.
\]
For the last equality, the invariant inner product makes the action unitary, so the trace of \(g^{-1}\) is the complex conjugate of the trace of \(g\). Equal characters therefore give equal multiplicities of every irreducible, and hence an equivariant isomorphism. Conversely, an equivariant isomorphism conjugates each pair of action matrices and preserves their traces. This proves the character criterion for finite-dimensional complex representations of a finite group.

\subsection{The seven exterior-power identities}\label{subsec:A2}

Use \((a,b,c,d)=a\mathbf1\oplus bA\oplus cB\oplus dC\), whose dimension is \(a+4b+5c+6d\). In every row \(\Lambda^2U\cong\Lambda^kV\). The columns \(5A,5B\) remain separate, even where their traces agree.

{\small
\begin{longtable}[]{@{}rllll@{}}
\toprule\noalign{}
Value & \((m,n;k)\) & \(U\) & \(V\) & Common character \\
\midrule\noalign{}
\endhead
\bottomrule\noalign{}
\endlastfoot
120 & \((16,10;3)\) & \((0,0,2,1)\) & \((1,1,1,0)\) & \((120,-8,3,0,0)\) \\
210 & \((21,10;4)\) & \((3,2,2,0)\) & \((1,1,1,0)\) & \((210,2,3,0,0)\) \\
1540 & \((56,22;3)\) & \((4,5,4,2)\) & \((3,2,1,1)\) & \((1540,-20,10,0,0)\) \\
3003 & \((78,15;5)\) & \((4,4,8,3)\) & \((0,1,1,1)\) & \((3003,-21,0,3,3)\) \\
7140 & \((120,36;3)\) & \((3,9,9,6)\) & \((1,1,5,1)\) & \((7140,-60,3,0,0)\) \\
11628 & \((153,19;5)\) & \((6,9,15,6)\) & \((1,0,0,3)\) & \((11628,-36,0,3,3)\) \\
24310 & \((221,17;8)\) & \((6,16,17,11)\) & \((3,1,2,0)\) & \((24310,-110,10,0,0)\) \\
\end{longtable}
}

The class sizes are \(1,15,20,12,12\). The characters of $\mathbf1,A,B,C$ are

{\small
\begin{longtable}[]{@{}lrrrrr@{}}
\toprule\noalign{}
Module & \(1\) & \(2A\) & \(3A\) & \(5A\) & \(5B\) \\
\midrule\noalign{}
\endhead
\bottomrule\noalign{}
\endlastfoot
\(1\) & 1 & 1 & 1 & 1 & 1 \\
\(A\) & 4 & 0 & 1 & -1 & -1 \\
\(B\) & 5 & 1 & -1 & 0 & 0 \\
\(C\) & 6 & -2 & 0 & 1 & 1 \\
\end{longtable}
}

The permutation and coset models determine the character on all sixty group elements. Both Newton recursion and determinant-polynomial multiplication give the character values displayed in the table. The characters of these modules take the same value on the two distinct conjugacy classes of order five.

\subsection{The three geometric comparisons}\label{subsec:A3}

For \(U=\mathbf1\oplus2E\), the character is \((21,5,3,1,1)\), so \(\chi_{\Lambda^2U}=(210,2,3,0,0)\). On the ten edges the nonidentity cycles are \(1^2 2^4\), \(1\,3^3\), and \(5^2\), giving exterior polynomials

\[
(1+t)^2(1-t^2)^4,\qquad (1+t)(1+t^3)^3,\qquad(1+t^5)^2.
\]

Their coefficients of \(t^4\) are \(2,3,0\). This proves the \(210\) identity. At a transposition in the corresponding \(S_5\) actions, the two characters are \(30\) and \(-14\); the identity concerns these restrictions to \(A_5\).

For \(U=\operatorname{Sym}^2Q\) and \(V=\Lambda^2P\), the symmetric- and exterior-square formulas give

\[
\chi_U=(78,6,0,3,3),\qquad \chi_V=(15,-1,0,0,0),
\qquad\chi_{\Lambda^2U}=(3003,-21,0,3,3).
\]

At an involution, \(V\) has seven eigenvalues \(+1\) and eight \(-1\), so its exterior polynomial is \((1-t)(1-t^2)^7\), whose coefficient of \(t^5\) is \(-21\). At order three each cube root of unity occurs with multiplicity five, giving \((1+t^3)^5\), whose coefficient of \(t^5\) is zero. At order five each fifth root of unity occurs with multiplicity three, giving \((1+t^5)^3\), whose coefficient of \(t^5\) is three. Together with the dimension, this proves equality on all five classes. Here $\dim V=15$; the equality with \(\binom{14}{6}\) is the classical Lind--Singmaster relation.

For \(W=Q\oplus B=\mathbf1\oplus S\) and \(T=\mathbf1\oplus3C\), the characters are \((17,1,-1,2,2)\) and \((19,-5,1,4,4)\). Thus \(\operatorname{Sym}^2W\) has character \((153,9,0,3,3)\), and \(\Lambda^2\operatorname{Sym}^2W\) has character \((11628,-36,0,3,3)\). The nonidentity exterior polynomials of \(T\) are

\[
(1+t)^7(1-t)^{12},\qquad(1+t)(1+t^3)^6,\qquad(1+t)^4(1+t^5)^3.
\]

Their coefficients of \(t^5\) are \(-36,0,3\), giving \(11628\). Together with the values at the identity, these character comparisons prove the isomorphisms of $A_5$-representations.

\subsection{Enumeration over $\mathbb C$ and $\mathbb Q$ and the exclusions at $3003$}\label{subsec:A4}

We count ordered pairs of complex representation isomorphism classes with the dimensions and exterior powers fixed. We order the irreducibles as $\mathbf1,D_+,D_-,A,B$, of dimensions $1,3,3,4,5$. With $\varphi=(1+\sqrt5)/2$, their characters are
\[
\begin{array}{c|rrrrr}
&1&2A&3A&5A&5B\\\hline
\mathbf1&1&1&1&1&1\\
D_+&3&-1&0&\varphi&1-\varphi\\
D_-&3&-1&0&1-\varphi&\varphi\\
A&4&0&1&-1&-1\\
B&5&1&-1&0&0
\end{array}
\]
Orthogonality with class sizes $(1,15,20,12,12)$ and the sum of the squares of the irreducible dimensions $1+9+9+16+25=60$ verify completeness. The rationally defined odd vertex module has complexification $C=D_+\oplus D_-$. A nonzero proper rational invariant subspace would complexify to one of $D_+$ or $D_-$. Neither constituent is Galois-stable, so the rational module is irreducible. The complexifications of rational representations are the nonnegative integer sums of $\mathbf1,A,B,C$; in complex coordinates they have equal multiplicities of $D_+,D_-$. All four summands have rational models.

Up to isomorphism, every representation of dimension $n$ is determined by a multiplicity tuple $(t,x,y,a,b)\in\mathbb Z_{\geq0}^5$ with
\[
t+3x+3y+4a+5b=n.
\]
Enumerating this finite set is exhaustive by complete reducibility. Its character is
\[
(n,\ t-x-y+b,\ t+a-b,\ t-a+y+(x-y)\varphi,\ t-a+x+(y-x)\varphi).
\]
Newton's identities compute exterior characters in $\mathbb Z[\varphi]$. The power maps are $2A^{2r}=1$, $3A^{3r}=1$, and for $5A$ the exponents $1,4$ stay in $5A$, $2,3$ lie in $5B$, and $0$ gives the identity; interchange $5A,5B$ for the other class. Matching the five exterior character values therefore decides isomorphism.

The resulting counts, after removing the single completely trivial pair, are
\[
\begin{array}{r|rr|rr}
\text{value}&\text{all complex pairs}&\text{rational pairs}&\#U\text{ candidates}&\#V\text{ candidates}\\\hline
120&10&4&73&22\\
210&2&2&159&22\\
1540&11&5&3853&181\\
3003&19&3&12614&62\\
7140&3&1&62029&856\\
11628&26&6&155358&118\\
24310&31&11&636215&85
\end{array}
\]
The candidate columns count all complex representations of the two prescribed dimensions. Distinct pairs related by exchanging $D_+$ and $D_-$ are counted separately; the two sides of a pair are ordered. The JSON certificate lists every surviving ordered pair and the numbers of candidate representations of the two prescribed dimensions. The verification program also checks that the pair printed in each article row belongs to the corresponding list.

For $210$, both pairs have $U=\mathbf1\oplus2E$ and $V\in\{E,A\oplus C\}$. Among representations defined over $\mathbb Q$, all three $3003$ pairs have $V=A\oplus B\oplus C$. The unique rational pair at $7140$, in the article's $(\mathbf1,A,B,C)$ coordinates, is $U=(3,9,9,6)$ and $V=(1,1,5,1)$.

The same exhaustive procedure on $(\dim U,\dim V;k)=(78,14;6)$ gives only $(78\mathbf1,14\mathbf1)$. For the different equation $\Lambda^5U\cong\Lambda^6V$ in dimensions $(15,14)$, matching the fifth and sixth exterior characters gives only $(15\mathbf1,14\mathbf1)$. The enumeration includes all complex representations in the specified dimensions. Together with the nontrivial $78/15$ construction, it proves that exactly one of the three pairwise exterior-power problems at $3003$ admits a nontrivial $A_5$-realization.

\subsection{The face construction and its uniqueness}\label{subsec:A5}

Let $G=A_5$. All representations in this section are complex, and a
permutation representation means one admitting a basis permuted by $G$.
We use the class order $(1,2A,3A,5A,5B)$, with class sizes
$(1,15,20,12,12)$. Let $X_5$ be the five-letter set, $X_{10}$ its unordered
pairs, and let $X_6,X_{12},X_{20}$ be the axes, vertices, and faces of the
icosahedron. Put $R_d=\mathbb C[X_d]$, with $P=R_6$, $E=R_{10}$ and
$Q=R_{12}$. The fixed identification of the icosahedral group with $A_5$
is used for all five actions.

\subsubsection*{Cosets, cycles, and the identity}

For an explicit model, the stabilizers of $X_5,X_6,X_{10},X_{12},X_{20}$
are respectively
\[
\begin{gathered}
\langle(123),(124)\rangle,\qquad
\langle(12345),(25)(34)\rangle,\qquad
\langle(345),(12)(34)\rangle,\\
\langle(12345)\rangle,\qquad \langle(123)\rangle.
\end{gathered}
\]
Their orders are $12,10,6,5,3$. Left multiplication on $G/H$ gives the
corresponding action. In particular the order-six subgroup stabilizes
$\{1,2\}$, identifying its cosets with the ten simplex edges. Face
stabilizers are generated by the rotations through $2\pi/3$ about their
centres and therefore give the last coset space.

For any subgroup $H$, counting fixed cosets gives
\[
\chi_{\mathbb C[G/H]}(g)
=\frac{|C_G(g)|}{|H|}\,|H\cap g^G|.
\]
The resulting fixed-label counts and cycle types are
\begin{center}\small
\begin{tabular}{c|rrr|lll}
\toprule
&\multicolumn{3}{c|}{Fixed labels}&\multicolumn{3}{c}{Cycle type}\\
Set&$2A$&$3A$&$5A,5B$&$2A$&$3A$&$5A,5B$\\\midrule
$X_5$&1&2&0&$1\,2^2$&$1^2 3$&$5$\\
$X_6$&2&0&1&$1^2 2^2$&$3^2$&$1\,5$\\
$X_{10}$&2&1&0&$1^2 2^4$&$1\,3^3$&$5^2$\\
$X_{12}$&0&0&2&$2^6$&$3^4$&$1^2 5^2$\\
$X_{20}$&0&2&0&$2^{10}$&$1^2 3^6$&$5^4$\\\bottomrule
\end{tabular}
\end{center}
Here $1^a r^b$ means $a$ fixed labels and $b$ cycles of length $r$.
For example an order-three rotation fixes the two faces whose centres lie
on its axis. A half-turn or a rotation of order five preserves no face.

\begin{proposition}\label{prop:faces-identity}
For $U=P\oplus E\oplus2R_{20}$ and $V=2R_5\oplus Q$, one has
\[
\dim U=56,\qquad\dim V=22,\qquad
\Lambda^2U\cong\Lambda^3V.
\]
The common character is $(1540,-20,10,0,0)$.
\end{proposition}

\begin{proof}
Adding the fixed-label counts gives
\[
\chi_U=(56,4,5,1,1),\qquad
\chi_V=(22,2,4,2,2).
\]
The formula $\chi_{\Lambda^2U}(g)
=(\chi_U(g)^2-\chi_U(g^2))/2$ gives the character in the statement.
The cycle types of $V$ at orders two, three, and five are
$1^2 2^{10}$, $1^4 3^6$, and $1^2 5^4$. A cycle of length $r$
contributes $1-(-t)^r$ to $\det(1+tg)$, so the cubic exterior characters are
\[
\begin{array}{c|c|r}
\text{class}&\det(1+tg|_V)&[t^3]\\\hline
1&(1+t)^{22}&1540\\
2A&(1+t)^2(1-t^2)^{10}&-20\\
3A&(1+t)^4(1+t^3)^6&10\\
5A&(1+t)^2(1+t^5)^4&0\\
5B&(1+t)^2(1+t^5)^4&0
\end{array}
\]
Thus the characters agree on all five classes, proving the isomorphism.
\end{proof}

In the article notation $(a,b,c,d)=a\mathbf1+bA+cB+dC$, the coset
decompositions are
\[
R_5=\mathbf1+A,\quad P=\mathbf1+B,\quad
E=\mathbf1+A+B,\quad Q=\mathbf1+B+C,\quad
R_{20}=\mathbf1+2A+B+C.
\]
Consequently $U=(4,5,4,2)$ and $V=(3,2,1,1)$. In coordinates relative to $(\mathbf1,D_+,D_-,A,B)$, these become
$(4,2,2,5,4)$ and $(3,1,1,2,1)$. The four orbits of the displayed
fifty-six-element set and the three orbits of the twenty-two-element set
give $\dim U^G=4$ and $\dim V^G=3$. Each invariant orbit sum is nonzero,
although neither displayed set has a singleton orbit.

\subsubsection*{A short uniqueness proof}

\begin{proposition}\label{prop:faces-unique}
Suppose $U'$ and $V'$ are permutation representations of $A_5$ of
dimensions $56$ and $22$ and
$\Lambda^2U'\cong\Lambda^3V'$. Then either both actions are trivial,
or $U'\cong P\oplus E\oplus2R_{20}$ and $V'\cong2R_5\oplus Q$.
\end{proposition}

We recall the complete transitive permutation catalogue. The group
$D_{10}$ below has order ten.
\begin{center}\small
\begin{tabular}{c|c|rrrrr}
\toprule
Module&Stabilizer&$1$&$2A$&$3A$&$5A$&$5B$\\\midrule
$R_1$&$G$&1&1&1&1&1\\
$R_5$&$A_4$&5&1&2&0&0\\
$R_6$&$D_{10}$&6&2&0&1&1\\
$R_{10}$&$S_3$&10&2&1&0&0\\
$R_{12}$&$C_5$&12&0&0&2&2\\
$R_{15}$&$V_4$&15&3&0&0&0\\
$R_{20}$&$C_3$&20&0&2&0&0\\
$R_{30}$&$C_2$&30&2&0&0&0\\
$R_{60}$&$1$&60&0&0&0&0\\\bottomrule
\end{tabular}
\end{center}
For completeness, a subgroup in the five-letter action either fixes a
letter, lies in the stabilizer $S_3$ of a $2+3$ partition, or is transitive.
The subgroups of $A_4$ and $S_3$ give all the intransitive entries above.
A transitive subgroup has one or six Sylow-five subgroups. With one it
lies in $D_{10}$; with six it contains all twenty-four five-cycles, so its
order is $30$ or $60$. Order $30$ would give an index-two quotient of
$A_5$, impossible because three-cycles generate $A_5$. This proves the
list. Within each type the subgroups are conjugate, as seen from the point
stabilizers and the cyclic subgroups and their normalizers. The fixed-coset
formula gives the displayed characters. Every permutation representation
is a nonnegative sum of these nine modules.

\begin{proof}[Proof of Proposition~\ref{prop:faces-unique}]
Let $f=\chi_{V'}(2A)$ and $s=\chi_{U'}(2A)$. They are nonnegative
fixed-label counts, and $f$ is even with $0\leq f\leq22$.
The exterior-cube formula at an involution gives
\[
\chi_{\Lambda^3V'}(2A)=\frac{f^3-3\cdot22f+2f}{6}
=\frac{f^3-64f}{6}.
\]
Equating this with $(s^2-56)/2$ yields
\[
s^2=\frac{f^3-64f+168}{3}.
\]
Checking the twelve possible even values of $f$ leaves exactly
\[
(f,s)=(2,4),\quad(6,0),\quad(22,56).
\]
\begin{samepage}
The twelve values are:
\begin{center}\small
\begin{tabular}{rr|rr}
\toprule
$f$&Required $s^2$&$f$&Required $s^2$\\\midrule
0&56&12&376\\
2&16&14&672\\
4&$-8$&16&1080\\
6&0&18&1616\\
8&56&20&2296\\
10&176&22&3136\\\bottomrule
\end{tabular}
\end{center}
\end{samepage}

If $(f,s)=(22,56)$, an involution fixes every basis label on both sides.
Its conjugates do so as well. Double transpositions generate $A_5$:
products such as $(12)(34)\,(12)(35)$ are three-cycles, and conjugates of
a three-cycle generate $A_5$. Thus both actions are trivial.

If $(f,s)=(6,0)$, every transitive summand of $U'$ has zero involution
trace. Since its dimension is $56$, the catalogue leaves only
$U'=aR_{12}\oplus bR_{20}$, with $12a+20b=56$. The only nonnegative
solution is $(a,b)=(3,1)$. Its order-three trace is two, giving exterior
square trace one. If an order-three element fixes $r$ labels of $V'$, its
other cycles have length three and
\[
\chi_{\Lambda^3V'}(3A)=\binom r3+\frac{22-r}{3},
\qquad r\in\{1,4,7,10,13,16,19,22\}.
\]
At $r=1$ this is seven, and for $r\geq4$ it is at least ten: its
increment when $r$ increases by three is $3r(r+1)/2>0$.
It therefore cannot equal one, excluding this case.

It remains to consider $(f,s)=(2,4)$. The dimension and involution trace
conditions on $V'$ leave just three orbit-sum presentations:
\[
V'=R_{10}\oplus R_{12},\qquad
V'=2R_1\oplus R_{20},\qquad
V'=2R_5\oplus R_{12}.
\]
Indeed only $R_1,R_5,R_6,R_{10}$ contribute to the allowed trace two;
$R_{12},R_{20}$ contribute zero. A copy of $R_6$ leaves sixteen dimensions
to be supplied by $R_{12},R_{20}$, which is impossible; the remaining
solutions of the two dimension equations are exactly those displayed.
The first presentation has order-three trace one and exterior-cube trace
seven. If $u=\chi_{U'}(3A)$, equality would require $u(u-1)/2=7$, whose
discriminant $57$ is not a square. The last two presentations have the
same character, namely $(22,2,4,2,2)$, and therefore define the same
representation $V$.

For this $V$, equality of exterior characters forces
$\chi_{U'}(3A)=5$, since $u(u-1)/2=10$ and $u\geq0$. The catalogue
also shows that the two order-five traces of any permutation module are
equal; write their common value as $q\geq0$. Its exterior-square trace
must be zero, giving $q(q-1)/2=0$ and $q\in\{0,1\}$.
The dimension of the invariant subspace of $U'$ is
\[
\frac{56+15\cdot4+20\cdot5+24q}{60}
=\frac{216+24q}{60}.
\]
Integrality forces $q=1$. Thus
$\chi_{U'}=(56,4,5,1,1)=\chi_U$, proving the assertion.
\end{proof}

There are several orbit presentations of this one representation pair.
For example,
\[
2R_5\oplus Q\cong2\mathbf1\oplus R_{20},
\]
although the orbit sizes of their basis sets are $(5,5,12)$ and
$(1,1,20)$. Similarly,
\[
P\oplus E\oplus2R_{20}
\cong\mathbf1\oplus R_{15}\oplus2R_{20}
\cong\mathbf1\oplus R_5\oplus R_{20}\oplus R_{30}.
\]
These are character equalities, not identifications of the corresponding
$G$-sets. They account for distinct presentations of the same ordered
representation pair.

\subsection{The orbits of vertex triples}\label{subsec:A6}

Let $G=A_5$ act by rotations of an icosahedron, with vertex set $X$.
All representations are complex. Write $Q=\mathbb C[X]$, let $P$ be
the permutation space on its six unoriented vertex axes, and let $R_5$
be the five-letter permutation space under a fixed identification of
the rotation group with $A_5$. We also use the face space $R_{20}$ and
the regular representation $R_{60}=\mathbb C[G]$. Conjugacy classes
are ordered $(1,2A,3A,5A,5B)$ and have sizes $(1,15,20,12,12)$.

\subsubsection*{Stabilizers and orientations}

\begin{proposition}\label{prop:triples-orbits}
The action on the $220$ three-element subsets of $X$ has orbit sizes
$20,20,60,60,60$. Its exterior basis can be reoriented to give a
permutation basis, and
\[
\Lambda^3Q\cong2R_{20}\oplus3R_{60}.
\]
\end{proposition}

\begin{proof}
Let $H$ be the setwise stabilizer of a triple $T$. Its action on $T$ is
faithful: an element in the kernel would fix three vertices, whereas
a nonidentity rotation fixes at most the two ends of its axis on the
sphere. Thus $H$ embeds in $S_3$. No involution preserves $T$, since
an involution acting on three points must fix one, and an icosahedral
half-turn fixes no vertex. Hence $H$ is trivial or cyclic of order
three. Its action on $T$ is even in either case.

There are twenty elements of order three in $G$, hence ten subgroups
of order three. Each such subgroup acts on $X$ in four cycles of
length three and preserves exactly those four triples. Each preserved
triple has that subgroup as its full stabilizer. Thus forty triples
have a stabilizer of order three; they form two orbits of size twenty.
The other $220-40=180$ triples have trivial stabilizers and form three
free orbits.

To pass from sets to exterior vectors, choose an ordering
$(x,y,z)$ of one representative $T$ in an orbit, and put
$w_T=e_x\wedge e_y\wedge e_z$. Every element of $H$ fixes $w_T$,
since its permutation of $T$ is even. The map
\[
\mathbb C[G/H]\longrightarrow
\operatorname{span}\{e_a\wedge e_b\wedge e_c:
                         \{a,b,c\}\in G\cdot T\},
\qquad [gH]\longmapsto g w_T
\]
is consequently well defined, equivariant, and an isomorphism: distinct
cosets have distinct underlying triples. All subgroups of order three
are conjugate, and a face stabilizer has order three. The two
twenty-element orbits therefore supply two copies of $R_{20}$, and
the three free orbits supply three copies of $R_{60}$.
\end{proof}

The five initial orientations may be reversed independently. These
choices specify the orbitwise identification of $\Lambda^3Q$; the
exterior identity below is established by character comparison.

\subsubsection*{The five kinds of triangles}

The orbit count has a geometric description. Scale the icosahedron to
have vertices
\[
(0,\pm1,\pm\varphi),\quad
(\pm1,\pm\varphi,0),\quad
(\pm\varphi,0,\pm1),\qquad
\varphi=\frac{1+\sqrt5}{2}.
\]
Distances between distinct vertices are the edge length $s=2$, the
long diagonal $\ell=2\varphi$, and the diameter
$d=2\sqrt{\varphi+2}$. The five orbits are
\begin{center}\small
\begin{tabular}{lccc}
\toprule
Triple type & Side lengths & Orbit size & Stabilizer\\\midrule
Ordinary face & $s,s,s$ &20&$C_3$\\
Large equilateral triangle & $\ell,\ell,\ell$ &20&$C_3$\\
Two edges & $s,s,\ell$ &60&$1$\\
One edge, no opposite pair & $s,\ell,\ell$ &60&$1$\\
An opposite pair and a third vertex & $s,\ell,d$ &60&$1$\\\bottomrule
\end{tabular}
\end{center}
Here is a direct count. In each of the graphs joining pairs at distance
$s$ and at distance $\ell$, every vertex has five neighbors, and the
graph induced on those neighbors is a pentagon; these assertions follow
by substitution in the displayed coordinates. Each graph consequently
has $12\cdot5/3=20$ triangles. There are also $6\cdot10=60$ triples
containing an opposite pair, with side lengths $s,\ell,d$.
If $a$ and $b$ count the remaining types $s,s,\ell$ and
$s,\ell,\ell$, then $a+b=120$. Counting incidences of a short edge
with a containing triple gives
\[
30\cdot10=3\cdot20+60+2a+b,
\]
so $a=b=60$. Side lengths are rotation invariants. By
Proposition~\ref{prop:triples-orbits}, each twenty-element family is
one orbit, and all remaining orbits are free. Each of the other three
families is therefore one orbit as well.

\subsubsection*{The exterior identity}

\begin{proposition}\label{prop:triples-identity}
For $U=\mathbf1\oplus\Lambda^3Q$ and $V=R_5\oplus2P$,
\[
\dim U=221,\qquad \dim V=17,\qquad
\Lambda^2U\cong\Lambda^8V.
\]
The common character is $(24310,-110,10,0,0)$.
\end{proposition}

\begin{proof}
At orders two, three, and five the vertex action has cycle types
$2^6$, $3^4$, and $1^2 5^2$. A cycle of length $r$ contributes
$1-(-t)^r$ to the exterior generating polynomial. Thus
\[
\chi_{\Lambda^3Q}=(220,0,4,0,0),\qquad
\chi_U=(221,1,5,1,1).
\]
The five-letter and axis characters are $(5,1,2,0,0)$ and
$(6,2,0,1,1)$, respectively, giving
$\chi_V=(17,5,2,2,2)$. The complete comparison is
\begin{center}\small
\begin{tabular}{c|c|r|r}
\toprule
Class&$\det(1+tg|_V)$&$[t^8]$&$\chi_{\Lambda^2U}$\\\midrule
$1$&$(1+t)^{17}$&24310&24310\\
$2A$&$(1+t)^5(1-t^2)^6$&$-110$&$-110$\\
$3A$&$(1+t)^2(1+t^3)^5$&10&10\\
$5A$&$(1+t)^2(1+t^5)^3$&0&0\\
$5B$&$(1+t)^2(1+t^5)^3$&0&0\\\bottomrule
\end{tabular}
\end{center}
The last column uses
$\chi_{\Lambda^2U}(g)=(\chi_U(g)^2-\chi_U(g^2))/2$.
At an involution the coefficient of $t^8$ is
$5\binom62-10\binom63+\binom64=75-200+15=-110$.
At order three only $t^2t^6$ contributes, giving $\binom52=10$;
at order five no exponent $5j+i=8$ has $0\leq i\leq2$.
The characters agree on all five classes.
\end{proof}

In the article notation $(a,b,c,d)=a\mathbf1+bA+cB+dC$, with
dimensions $(1,4,5,6)$, one has
\[
\Lambda^3Q=(5,16,17,11),\qquad
U=(6,16,17,11),\qquad V=(3,1,2,0).
\]
In coordinates relative to $(\mathbf1,D_+,D_-,A,B)$, the three tuples are
\[
(5,11,11,16,17),\qquad (6,11,11,16,17),\qquad (3,0,0,1,2).
\]
Here $C=D_+\oplus D_-$, $R_5=\mathbf1\oplus A$, and
$P=\mathbf1\oplus B$. The orbit decomposition gives
$\dim U^G=6$, while the three orbits defining $V$ give $\dim V^G=3$.

\Needspace{22\baselineskip}
\subsubsection*{A second seventeen-dimensional representation}

Let $B$ be the zero-sum subspace of $P$ and put
$W=\mathbf1\oplus3B$, of dimension sixteen. The module
$\widetilde V=\mathbf1\oplus W=2\mathbf1\oplus3B$ has dimension
seventeen and is not isomorphic to $V=R_5\oplus2P$.

\begin{proposition}\label{prop:triples-alternative}
One has $V\not\cong\widetilde V$ and
$\Lambda^8V\cong\Lambda^8\widetilde V$.
\end{proposition}

\begin{proof}
Their characters are
\[
\chi_V=(17,5,2,2,2),\qquad
\chi_{\widetilde V}=(17,5,-1,2,2),
\]
so the order-three traces distinguish them. Since
$P=\mathbf1\oplus B$, its determinant polynomial gives
\[
\det(1+tg|_{\widetilde V})
=\frac{\det(1+tg|_P)^3}{1+t}.
\]
At orders two and five this agrees with the polynomial for $V$ in
the preceding table. At order three it is
\[
\frac{(1+t^3)^6}{1+t}=(1+t)^5(1-t+t^2)^6,
\]
whose coefficient of $t^8$ is $1-6+15=10$. At the identity both eighth
exterior powers have dimension $\binom{17}{8}$. Their characters
therefore agree on all classes.
\end{proof}

Consequently,
$\Lambda^2(\mathbf1\oplus\Lambda^3Q)
\cong\Lambda^8(\mathbf1\oplus W)$.
Every permutation trace counts fixed labels, so the order-three value
$-1$ rules out a permutation basis for $\widetilde V$.

\subsection{Cubic stabilizer characters}\label{subsec:A7}

Let $G=A_5$. All representations are complex. Let $R_5$ and $P$ be
the permutation spaces on five letters and six icosahedral axes, and
write $R_5=\mathbf1\oplus A$, $P=\mathbf1\oplus B$.
In class order $(1,2A,3A,5A,5B)$, with sizes $(1,15,20,12,12)$,
\[
\chi_A=(4,0,1,-1,-1),\qquad \chi_B=(5,1,-1,0,0).
\]
These follow by subtracting the constant functions from the fixed-label
characters. Let $R_{20}$ be the face permutation space, whose character
is $(20,0,2,0,0)$, and set $\omega=e^{2\pi i/3}$.

\subsubsection*{Lines over an orbit}

For a subgroup $H\leq G$ and a character
$\theta:H\to\mathbb C^\times$, the induced representation
\[
L(H,\theta)=\mathbb C[G]\otimes_{\mathbb C[H]}\mathbb C_\theta
\]
is a direct sum of one-dimensional spaces indexed by $G/H$.
Choose coset representatives $r_i$ and write $e_i=r_i\otimes1$.
If $g r_i=r_jh$, then
\[
g e_i=\theta(h)e_j.
\]
Thus a cubic character gives a basis whose nonzero matrix entries
belong to $\{1,\omega,\omega^2\}$. Tracing the fixed cosets gives
\begin{equation}\label{eq:lines-induced-character}
\chi_{L(H,\theta)}(g)
=\frac{|C_G(g)|}{|H|}\sum_{h\in H\cap g^G}\theta(h).
\end{equation}
Indeed $xH$ is fixed exactly when $x^{-1}gx\in H$, in which case its
line contributes $\theta(x^{-1}gx)$. Each conjugate in $H$ is obtained
from $|C_G(g)|$ choices of $x$, and a coset has $|H|$ representatives.

\begin{proposition}\label{prop:lines-inductions}
Let $H=A_4$ be a letter stabilizer and let $\eta$ be either nontrivial
character of $H/V_4$. If $K=C_3$ is a face stabilizer and $\theta$ is
either nontrivial character of $K$, then
\[
L(H,\eta)\cong B,\qquad L(K,\theta)\cong T:=A\otimes B.
\]
\end{proposition}

\begin{proof}
The three partitions of the other four letters into two pairs are
permuted cyclically by $A_4$. The kernel is its Klein four subgroup,
so this realizes $H/V_4\cong C_3$ and constructs $\eta$.
The intersections of $H$ with the five $G$-classes have character sums
\[
(1,3,4\omega+4\omega^2,0,0)=(1,3,-4,0,0).
\]
The centralizer orders are $(60,4,3,5,5)$.
Formula~\eqref{eq:lines-induced-character}, with $|H|=12$, gives
$(5,1,-1,0,0)=\chi_B$.

For $K$, the character sums are
$(1,0,\omega+\omega^2,0,0)=(1,0,-1,0,0)$.
Dividing the corresponding centralizer multiples by three gives
\[
\chi_{L(K,\theta)}=(20,0,-1,0,0)
=\chi_A\chi_B=\chi_{A\otimes B}.
\]
Equality of characters proves both isomorphisms. Replacing either
cubic character by its inverse leaves its induced character unchanged.
\end{proof}

The characters belong to the stabilizers of the lines. Every one-dimensional character of $G$ is trivial: a three-cycle is conjugate in
$A_5$ to its inverse, so its image under a linear character has both
square and cube equal to one. Three-cycles generate $A_5$.

\subsubsection*{The tangent lines at face centres}

Let $F$ be the twenty face centres projected radially onto the unit
sphere. Give the real tangent plane at $p\in F$ the complex structure
\[
J_p(v)=p\times v,\qquad v\perp p.
\]
Since $J_p^2=-1$, $J_p$ makes the real plane $T_pS^2$ into a complex line. Every rotation $g$ satisfies
$g(p\times v)=(gp)\times(gv)$, so its tangent map is complex linear.
The direct sum
\[
\mathcal T=\bigoplus_{p\in F}(T_pS^2,J_p)
\]
is therefore a twenty-dimensional complex representation. At a chosen
face centre, the stabilizer acts through a nontrivial cubic character.
Transport of this one line around the orbit identifies
$\mathcal T\cong L(C_3,\theta)\cong T$.

Its character also has a direct geometric calculation. An order-three
rotation fixes two face centres at opposite ends of its axis. Their
outward orientations give the scalars $\omega$ and $\omega^2$, whose
sum is $-1$. Rotations of orders two and five fix no face centre and
give trace zero. At the identity the trace is twenty, hence
$\chi_{\mathcal T}=(20,0,-1,0,0)$.

\subsubsection*{The identity and bases with cube roots of unity}

\begin{proposition}\label{prop:lines-identity}
Set $W=\mathbf1\oplus3B$, $U=\Lambda^2W$, and $V=W\oplus T$.
Then
\[
\dim W=16,\quad\dim U=120,\quad\dim V=36,\qquad
\Lambda^2U\cong\Lambda^3V.
\]
Moreover $U\cong3R_{20}\oplus3T$, so each of $W,U,V$ admits a basis
whose nonzero action coefficients are cube roots of unity.
\end{proposition}

\begin{proof}
The characters of $B,T,W,U,V$ are
\begin{center}\small
\begin{tabular}{c|rrrrr}
\toprule
& $1$&$2A$&$3A$&$5A$&$5B$\\\midrule
$B$&5&1&$-1$&0&0\\
$T$&20&0&$-1$&0&0\\
$W=\mathbf1\oplus3B$&16&4&$-2$&1&1\\
$U=\Lambda^2W$&120&0&3&0&0\\
$V=W\oplus T$&36&4&$-3$&1&1\\\bottomrule
\end{tabular}
\end{center}
Here $\chi_{\Lambda^2W}(g)
=(\chi_W(g)^2-\chi_W(g^2))/2$. In particular
\[
\chi_U=3(\chi_{R_{20}}+\chi_T),\qquad
U\cong3R_{20}\oplus3T.
\]
This decomposition and Proposition~\ref{prop:lines-inductions} give bases with coefficients in $\{1,\omega,\omega^2\}$:
\begin{align*}
W&\cong\mathbf1\oplus3L(A_4,\eta),\\
U&\cong3L(C_3,1)\oplus3L(C_3,\theta),\\
V&\cong\mathbf1\oplus3L(A_4,\eta)\oplus L(C_3,\theta).
\end{align*}
The usual wedge basis of $\Lambda^2W$ includes ordering signs. The induced decomposition above gives a basis of $U$ with coefficients in $\{1,\omega,\omega^2\}$.

\begin{samepage}
For the exterior identity, the determinant polynomials of $V$ and the
exterior-square traces of $U$ give the complete comparison
\begin{center}\small
\begin{tabular}{c|c|r|r}
\toprule
Class&$\det(1+tg|_V)$&$[t^3]$&$\chi_{\Lambda^2U}$\\\midrule
$1$&$(1+t)^{36}$&7140&7140\\
$2A$&$(1+t)^4(1-t^2)^{16}$&$-60$&$-60$\\
$3A$&$(1+t)^{10}(1-t+t^2)^{13}$&3&3\\
$5A$&$(1+t)(1+t^5)^7$&0&0\\
$5B$&$(1+t)(1+t^5)^7$&0&0\\\bottomrule
\end{tabular}
\end{center}
\end{samepage}
At order two the eigenvalue multiplicities are $20$ and $16$, giving
the coefficient of $t^3$, namely $4-64=-60$. At order three they are $10,13,13$
for $1,\omega,\omega^2$. Thus its polynomial can also be written
$(1+t^3)^{13}/(1+t)^3$, whose coefficient of $t^3$ is $13-10=3$.
At order five the multiplicities are eight for $1$ and seven for each
nontrivial fifth root, so the coefficient of $t^3$ is zero. These
multiplicities follow from the character of $V$ and its values on the
powers of the element. At the identity the coefficient of $t^3$ is
$\binom{36}{3}=7140$.

On $U$, the square formula gives
$\binom{120}{2}=7140$, $(0^2-120)/2=-60$, $(3^2-3)/2=3$, and
zero on both order-five classes. This proves the isomorphism.
\end{proof}

In article coordinates $(a,b,c,d)=a\mathbf1+bA+cB+dC$, where
$C=D_+\oplus D_-$, the modules are
\[
W=(1,0,3,0),\quad T=(0,1,2,1),\quad
U=(3,9,9,6),\quad V=(1,1,5,1).
\]
Here $(\dim A,\dim B,\dim C)=(4,5,6)$.

\subsubsection*{Why signs are insufficient}

\begin{proposition}\label{prop:lines-signs}
A signed permutation representation has nonnegative trace at every
element of order three. Consequently $V=W\oplus T$ admits no signed
permutation basis.
\end{proposition}

\begin{proof}
The permutation of a signed basis induced by $g$, with $g^3=1$, has
cycles of length one or three. A cycle of length three contributes
zero to the trace. On a fixed vector the coefficient is
$\epsilon\in\{1,-1\}$ and $\epsilon^3=1$, hence $\epsilon=1$.
The trace is the number of fixed basis vectors. Since $\chi_V(3A)=-3$,
no such basis exists for $V$.
\end{proof}

\subsubsection*{The exterior-square root}

\begin{proposition}\label{prop:lines-root}
Up to isomorphism, $W=\mathbf1\oplus3B$ is the unique
sixteen-dimensional complex $A_5$-module whose exterior square is isomorphic to $U$, where $\chi_U=(120,0,3,0,0)$.
\end{proposition}

\begin{proof}
Every element of $A_5$ is conjugate to its inverse, so every complex
character of $A_5$ is real-valued. Write the character of a possible
root as $(16,s,r,x,y)$. The square formula gives
\[
s^2=16,\qquad r^2-r=6,\qquad x^2=y,\qquad y^2=x.
\]
Thus $s=\pm4$, $r\in\{3,-2\}$, and $x=y=q\in\{0,1\}$.
For the last assertion, $x^4=x$ and $x$ is real.
The character inner product gives the multiplicity of $A$ as
$(64+20r-24q)/60$.
For $(r,q)=(3,0),(3,1),(-2,0),(-2,1)$, the numerators are
$124,100,24,0$. Integrality forces $r=-2$ and $q=1$.
The multiplicity of the trivial representation is then
\[
\frac{16+15s-40+24}{60}=\frac{s}{4}.
\]
Nonnegativity forces $s=4$. The character is therefore
$(16,4,-2,1,1)=\chi_{\mathbf1\oplus3B}$, which determines the
module up to isomorphism. Proposition~\ref{prop:lines-identity} gives the representation with this character.
\end{proof}

\subsection{Permutation, signed, and monomial counts}\label{subsec:A8}

A representation is a \emph{permutation representation} if it has a basis permuted by $G$. It is a \emph{signed permutation representation} if some basis is permuted with coefficients in $\{1,-1\}$. It is a \emph{complex monomial representation} if some basis is permuted with coefficients in $\mathbb C^\times$. These are properties of the representation isomorphism class; the basis is allowed to change.

\begin{proposition}[Induced line description]\label{l:monomial}
A complex monomial representation is exactly a direct sum of representations $\operatorname{Ind}_H^G\theta$ with $\theta:H\to\mathbb C^\times$ a linear character. The permutation and signed classes are obtained by requiring, respectively, $\theta=1$ or $\theta(H)\subseteq\{1,-1\}$.
\end{proposition}
\begin{proof}
Decompose a monomial basis into orbits of its one-dimensional spans. Choose a vector $e$ in one such line. Its stabilizer $H$ acts as $he=\theta(h)e$, and the action law makes $\theta$ a homomorphism. The orbit spans a copy of $\mathbb C[G]\otimes_{\mathbb C[H]}\mathbb C_\theta$: the map $g\otimes1\mapsto ge$ is well defined and sends coset basis vectors to independent lines. Taking every orbit gives the asserted direct sum. Conversely, choose coset representatives $r_i$. If $g r_i=r_jh$, the induced action sends $r_i\otimes1$ to $\theta(h)r_j\otimes1$, which is monomial. This proves the two restricted statements as well. As $H$ is finite, every value of $\theta$ is a root of unity.
\end{proof}

A geometric version consists of a finite $G$-set with a one-dimensional complex vector space above each point and a compatible linear action between fibers. Its direct sum of fibers is the representation in Proposition~\ref{l:monomial}. On a transitive orbit, the stabilizer character specifies this construction.

The nine transitive permutation characters and the subgroup completeness argument appear in Section~\ref{subsec:A5}. Their stabilizers have the following counts and abelianizations.
\begin{center}
\begin{tabular}{c|rrrrrrrrr}
$H$&$1$&$C_2$&$C_3$&$V_4$&$C_5$&$S_3$&$D_{10}$&$A_4$&$G$\\\hline
Number&1&15&10&5&6&10&6&5&1\\
$H^{\rm ab}$&$1$&$C_2$&$C_3$&$V_4$&$C_5$&$C_2$&$C_2$&$C_3$&$1$
\end{tabular}
\end{center}
The linear characters of each subgroup $H$ are precisely the characters of its abelianization $H^{\mathrm{ab}}$. For example $[A_4,A_4]=V_4$ and $[D_{10},D_{10}]=C_5$; for $A_5$ the commutator subgroup contains a three-cycle and hence all three-cycles, which generate it. Summing $|H^{\rm ab}|$ over all $59$ subgroups gives $159$ linear characters.

The nontrivial linear characters give eight further distinct representations:
\begin{center}\small
\begin{tabular}{cll}
\toprule
$H$&Character&Induced representation\\\midrule
$D_{10}$&nontrivial sign&$C$\\
$S_3$&sign&$A+C$\\
$V_4$&any of the three nontrivial characters&$A+B+C$\\
$C_2$&nontrivial character&$2A+2B+2C$\\
$A_4$&either nontrivial cubic character&$B$\\
$C_3$&either nontrivial cubic character&$A+2B+C$\\
$C_5$&$h\mapsto\zeta^{\pm1}$&$D_++A+B$\\
$C_5$&$h\mapsto\zeta^{\pm2}$&$D_-+A+B$\\\bottomrule
\end{tabular}\end{center}
Here $h\in5A$, $\zeta=e^{2\pi i/5}$, and the labels are chosen so that $D_+$ has eigenvalues $1,\zeta,\zeta^{-1}$ at $h$. These induced characters follow directly from~\eqref{eq:lines-induced-character}. For instance the nontrivial $V_4$ character has class sums $(1,-1,0,0,0)$ and induces $(15,-1,0,0,0)$; the sign on $D_{10}$ has sums $(1,-5,0,2,2)$ and induces $(6,-2,0,1,1)$. For $C_5$, Frobenius reciprocity reads off the multiplicity of $\zeta^j$ in the displayed eigenvalue lists, giving the last two entries. These calculations account for every character in each abelianization.

\begin{proposition}[Three finite catalogues]\label{p:catalogue}
The permutation representations are exactly the nonnegative sums of the nine transitive modules in Section~\ref{subsec:A5}. The signed permutation representations are exactly the nonnegative sums of those nine and the first four entries of the table above. The complex monomial representations are exactly the nonnegative sums of all seventeen entries.
\end{proposition}
\begin{proof}
Every orbit of basis lines is induced from a stabilizer character by Proposition~\ref{l:monomial}. The subgroup classification and the two tables list every such induced character. Conversely each listed summand has the required basis, and direct sums preserve that property.
\end{proof}

\subsubsection*{Exact counts and their completeness}
For fixed dimensions $(m,n)$ and degree $k$, we count ordered pairs of representation isomorphism classes $(U,V)$ satisfying $\Lambda^2U\cong\Lambda^kV$, excluding $(m\mathbf1,n\mathbf1)$. The rational column counts complexifications of representations defined over $\mathbb Q$. By Section~\ref{subsec:A4}, these are the nonnegative sums of $\mathbf1,A,B,C$, characterized by equal multiplicities of $D_+$ and $D_-$ in each representation.

\begin{proposition}[Enumeration]\label{t:counts}
The counts in the seven prescribed dimension triples are as follows. Each restricted column requires both $U$ and $V$ to belong to that class.
\begin{center}\small
\begin{tabular}{rcrrrrr}
\toprule
Value&$(m,n;k)$&Complex&Rational&Perm.&Signed&Monomial\\\midrule
120&$(16,10;3)$&10&4&0&1&3\\
210&$(21,10;4)$&2&2&1&2&2\\
1540&$(56,22;3)$&11&5&1&3&11\\
3003&$(78,15;5)$&19&3&0&2&3\\
7140&$(120,36;3)$&3&1&0&0&1\\
11628&$(153,19;5)$&26&6&0&2&8\\
24310&$(221,17;8)$&31&11&2&4&15\\\bottomrule
\end{tabular}\end{center}
In particular, the pair in Proposition~\ref{prop:faces-identity} is the unique nontrivial permutation pair at $1540$, up to ordered representation isomorphism.
\end{proposition}
\begin{proof}
Every complex module of dimension $d$ has a unique tuple
\[
(t,x,y,a,b)\in\mathbb Z_{\geq0}^5,\qquad t+3x+3y+4a+5b=d.
\]
Enumerate the nonnegative quadruples $(x,y,a,b)$ satisfying $3x+3y+4a+5b\leq d$, and set $t=d-3x-3y-4a-5b$. This gives every multiplicity tuple exactly once. The numbers of candidates are
\begin{center}\small
\begin{tabular}{r|rrrrrrr}
Value&120&210&1540&3003&7140&11628&24310\\\hline
$U$&73&159&3853&12614&62029&155358&636215\\
$V$&22&22&181&62&856&118&85
\end{tabular}\end{center}

For every candidate $U$ and $V$, compute $\chi_{\Lambda^2U}$ and $\chi_{\Lambda^kV}$. Pair $U$ with precisely those $V$ for which these characters agree. Removing the completely trivial pair leaves respectively $10,2,11,19,3,26,31$ pairs. The complete list of $102$ pairs, their common characters, and their membership in the three restricted classes is given in the computational supplement.

Here is a precise membership test. Order one of the generator lists as $g_1,\ldots,g_s\in\mathbb Z_{\geq0}^5$. For a target $v$, success at the first index means that there are integers $a_i\geq0$ with $v=\sum a_i g_i$. At index $i$, test each integer
\[
0\leq a_i\leq\min_{j:(g_i)_j>0}
                     \left\lfloor\frac{v_j}{(g_i)_j}\right\rfloor
\]
and recur on $v-a_i g_i$ and the next index. The recursion finds a decomposition when the residual vector is zero. A nonzero residual after the final generator gives no decomposition. Every branch is finite, and every possible expression satisfies these bounds. Induction on the number of remaining generators proves the correctness of both success and failure. The returned coefficient vector is recorded whenever the test succeeds. Each representation pair is counted once, regardless of how many decompositions into induced representations the membership test finds. Proposition~\ref{p:catalogue} proves that these tests classify exactly the three representation classes.

All arithmetic is integral. \texttt{verify\_revision.py} uses Newton's recurrence in $\mathbb Z[\varphi]$, with exact division at every step. \texttt{induction/verify\_wave\_i.py} uses eigenvalue factors instead. For an element of order $r$, encode an eigenvalue as $z^a$ in $\mathbb Z[z]/(z^r-1)$ and multiply $1+t z^a$, retaining coefficients through the needed exterior degree. The exponent lists for the nonidentity orders are
\begin{center}\small
\begin{tabular}{c|ccc}
&$r=2$&$r=3$&$r=5$ in class $5A$\\\hline
$\mathbf1$&$0$&$0$&$0$\\
$D_+$&$0,1,1$&$0,1,2$&$0,1,4$\\
$D_-$&$0,1,1$&$0,1,2$&$0,2,3$\\
$A$&$0,0,1,1$&$0,0,1,2$&$1,2,3,4$\\
$B$&$0,0,0,1,1$&$0,1,1,2,2$&$0,1,2,3,4$
\end{tabular}\end{center}
For $5B$ double each exponent modulo five. These lists follow by taking traces of all powers in the character table and Fourier inversion on the cyclic subgroup. After multiplication, a symmetric coefficient vector $(a_0,\ldots,a_{r-1})$ reduces to
\[
a_0-a_1\quad(r=2,3),\qquad
(a_0-a_1)+(a_1-a_2)\varphi\quad(r=5).
\]
The symmetry conditions are $a_1=a_2$ for $r=3$ and $a_1=a_4,a_2=a_3$ for $r=5$; the program checks them. This computation and Newton's recurrence give the same five character values, using exact arithmetic.

For an additional check of the catalogue, the program constructs $G$ as all even permutations, closes generated subgroups under multiplication, and repeatedly adjoins each missing element. Any subgroup can be reached by successively adjoining its generators, proving completeness of this construction. For every $H$ it checks every normal subgroup $K\leq H$. When $H/K$ is cyclic, each choice of a primitive root for its generator defines one faithful character of the quotient. These are all linear characters of $H$, since every finite subgroup of $\mathbb C^\times$ is cyclic and $K$ is the kernel. The code does not restrict the possible quotient orders in advance. It finds only $1,2,3,5$, checks every character's multiplication law, and checks that their number is $|H/[H,H]|$. Induction gives exactly the $9$, $13$, and $17$ distinct generator tuples above.

The two independent exterior-character computations give the same complete lists of ordered pairs. The catalogue membership test then gives the three restricted counts.
\end{proof}

\begin{proposition}[The signed exclusion at $7140$]\label{p:7140all}
None of the three nontrivial complex pairs in dimensions $(120,36;3)$ consists of two signed permutation representations. Exactly one pair consists of two complex monomial representations.
\end{proposition}
\begin{proof}
The complete list, in five-coordinate notation, is
\[
\begin{array}{c|c}
U&V\\\hline
(2,5,6,10,9)&(1,2,0,1,5)\\
(2,6,5,10,9)&(1,0,2,1,5)\\
(3,6,6,9,9)&(1,1,1,1,5).
\end{array}
\]
Every $V$ has trace $t+a-b=1+1-5=-3$ at a three-cycle. Proposition~\ref{prop:lines-signs} excludes a signed basis for each of them. The last pair has the bases with cube roots of unity constructed in Proposition~\ref{prop:lines-identity}. To see directly why the other two are not monomial, use the catalogue: the only generators with unequal $D_+,D_-$ multiplicities are $D_++A+B$ and $D_-+A+B$. A difference of two requires at least two such summands and hence at least two copies of $A$. Each of the first two $V$'s has only one copy of $A$. This proves the exclusion.
\end{proof}

The two permutation pairs at $24310$, in article coordinates, are
\[
((6,16,17,11),(3,1,2,0)),\qquad
((12,22,11,11),(5,3,0,0)).
\]
The second is $U=\mathbf1\oplus11R_{20}$ and $V=2\mathbf1\oplus3R_5$. Different orbit-sum expressions of these modules give the same two representation pairs.

\subsubsection*{Other realizations at $1540$ and $24310$}
The following pairs also satisfy the identities at $1540$ and $24310$. In coordinates relative to $(\mathbf1,A,B,C)$, they are
\[
\begin{array}{r|c|c|c}
\text{value}&U&V&\text{common exterior character}\\\hline
1540&(0,2,6,3)&(4,2,2,0)&(1540,-28,10,0,0)\\
24310&(3,13,20,11)&(2,0,3,0)&(24310,-110,10,0,0).
\end{array}
\]
At $1540$ these representations have characters $(56,0,-4,1,1)$ and
$(22,6,4,2,2)$; Newton's recurrence gives the displayed exterior character.
At $24310$ their characters are $(221,1,-4,1,1)$ and $(17,5,-1,2,2)$.
The latter $V$ is the alternative representation of Proposition~\ref{prop:triples-alternative}.
The article's constructions at $1540$ and $24310$ have permutation models described
in Sections~\ref{subsec:A5} and~\ref{subsec:A6}.

\section{Signed orbits and basis matching}\label{sec:orbits}

The signed bases of \(\Lambda^2\operatorname{Sym}^2Q\) and \(\Lambda^5\Lambda^2P\) have different unsigned orbit distributions: the left has \(39,20,3,1,1\) orbits of sizes \(60,30,15,12,6\), and the right has \(40,18,3,3\) of sizes \(60,30,15,6\). Therefore there is no equivariant bijection of these two particular unsigned bases.

This is compatible with their representation isomorphism, whose signed-orbit comparison reduces, after cancelling the common direct summands, to

\[
\operatorname{Ind}_{C_2}^{G}(\mathbf1\oplus\varepsilon)\cong\mathbb C[G],\qquad
\operatorname{Ind}_{C_5}^{G}\mathbf1\cong
\operatorname{Ind}_{D_5}^{G}(\mathbf1\oplus\varepsilon).
\]

Here \(D_5\) has order ten and its sign character factors through \(D_5/C_5\). The first equality is the induction of a regular representation; the second is induction in stages. The geometric construction in Section F gives a separate, factored linear intertwiner.

The full signed stabilizer enumeration is as follows; \(D_5\) has order ten.

{\small
\begin{longtable}[]{@{}lrr@{}}
\toprule\noalign{}
Stabilizer and sign character & Left orbits & Right orbits \\
\midrule\noalign{}
\endhead
\bottomrule\noalign{}
\endlastfoot
\texttt{1/trivial} & 39 & 40 \\
\texttt{C2/nontrivial} & 14 & 13 \\
\texttt{C2/trivial} & 6 & 5 \\
\texttt{V4/nontrivial} & 3 & 3 \\
\texttt{C5/trivial} & 1 & 0 \\
\texttt{D5/nontrivial} & 1 & 2 \\
\texttt{D5/trivial} & 0 & 1 \\
\end{longtable}
}

The common summand has dimension \(2931\). The remaining summands have dimensions \(30+30+12\) and \(60+6+6\) and are identified by the induction formulas above. Each label denotes the representation induced from the indicated one-dimensional character of its stabilizer. The finite orbit calculation also identifies the conjugacies among these signed stabilizers. For the Klein four group, its normalizer permutes its three nontrivial one-dimensional characters.

\section{Obstruction for complex semisimple Lie algebras}\label{sec:obstruction}

\begin{proposition}[Semisimple obstruction]\label{prop:C1}
Let $\mathfrak g$ be a finite-dimensional complex semisimple Lie algebra, and let $U,V$ be finite-dimensional complex $\mathfrak g$-modules, possibly reducible. If $\Lambda^2U\cong\Lambda^kV$ and
\[
\begin{gathered}
(\dim U,\dim V,k)\in
\{(21,10,4),(56,22,3),(78,14,6),(78,15,5),\\
(120,36,3),(153,19,5),(221,17,8)\},
\end{gathered}
\]
then both actions are trivial.
\end{proposition}

In particular, $\mathfrak g=\mathfrak g_2$, $U=\mathfrak e_6|_{\mathfrak g_2}$ and $V=\mathfrak g_2$ give a negative answer to the proposed $78/14$ exterior-power isomorphism in Baez's question \citep{Baez2026Coincidences}.

Restrict to a root copy of \(\mathfrak{sl}_2\). Its irreducible representation \(R_d\) of dimension \(d\) has weights \(d-1,d-3,\ldots,1-d\). Every \(\mathfrak{sl}_2\)-module of dimension \(n\) is a direct sum of irreducibles whose dimensions form a partition of \(n\). Thus the partitions give the complete finite list of possibilities. Put

\[
S_r(W)=\sum_{w\text{ a weight of }W} w^r.
\]

Weights are integral and occur in opposite pairs. Therefore \(S_2\) is an even nonnegative integer, vanishing exactly on a trivial module, and

\[
S_4-S_2\in24\mathbb Z,\qquad S_6-S_2\in120\mathbb Z.
\]

Indeed, \(a^4-a^2\) is divisible by twelve and \(a^6-a^2\) by sixty for every integer \(a\), and paired weights double these contributions. Expanding sums of distinct pairs of weights gives

\[
\begin{aligned}
S_2(\Lambda^2U)&=(m-2)S_2(U),\\
S_4(\Lambda^2U)&=(m-8)S_4(U)+3S_2(U)^2,\\
S_6(\Lambda^2U)&=(m-32)S_6(U)+15S_2(U)S_4(U).
\end{aligned}
\]

These follow equivalently by differentiating

\[
f_{\Lambda^2U}(t)=\frac{f_U(t)^2-f_U(2t)}2,\qquad f_U(t)=\sum_we^{tw},
\]

at orders two, four and six, using the vanishing of odd moments.

Given a partition defining \(V\), the moments of \(\Lambda^kV\) force candidate values of the three moments of \(U\) through these formulas. The denominators \(m-2,m-8,m-32\) are nonzero in every listed case. Necessary conditions are positivity and integrality for a nontrivial module, evenness of \(S_2\), and the two congruences above. To exclude every nontrivial candidate for \(V\), it suffices to show that the forced moments of \(U\) violate at least one of these necessary conditions.

{\small
\begin{longtable}[]{@{}
  >{\raggedright\arraybackslash}p{(\columnwidth - 8\tabcolsep) * \real{0.1579}}
  >{\raggedleft\arraybackslash}p{(\columnwidth - 8\tabcolsep) * \real{0.2105}}
  >{\raggedleft\arraybackslash}p{(\columnwidth - 8\tabcolsep) * \real{0.2105}}
  >{\raggedleft\arraybackslash}p{(\columnwidth - 8\tabcolsep) * \real{0.2105}}
  >{\raggedleft\arraybackslash}p{(\columnwidth - 8\tabcolsep) * \real{0.2105}}@{}}
\toprule\noalign{}
\((m,n,k)\) & Nontrivial partitions & Satisfy second-moment conditions & Satisfy fourth-moment conditions & Satisfy sixth-moment conditions \\
\midrule\noalign{}
\endhead
\bottomrule\noalign{}
\endlastfoot
\((21,10,4)\) & 41 & 1 & 0 & 0 \\
\((56,22,3)\) & 1001 & 25 & 1 & 0 \\
\((78,14,6)\) & 134 & 5 & 0 & 0 \\
\((78,15,5)\) & 175 & 2 & 0 & 0 \\
\((120,36,3)\) & 17976 & 161 & 1 & 0 \\
\((153,19,5)\) & 489 & 0 & 0 & 0 \\
\((221,17,8)\) & 296 & 1 & 0 & 0 \\
\end{longtable}
}

The two candidates remaining after the fourth-moment test are \(V=R_7\oplus R_5\oplus5R_2\) in dimension twenty-two and \(V=R_9\oplus R_8\oplus R_5\oplus2R_3\oplus4R_2\) in dimension thirty-six. The former forces \((S_2,S_4,S_6)=(570,15522,557040)\), with \(S_6-S_2\equiv30\pmod{120}\). The latter forces \((S_2,S_4,S_6)=(2244,115908,95867424/11)\), whose last entry is nonintegral.

Thus \(V\) is trivial on every root subalgebra, and the second-moment identity forces the same conclusion for \(U\). Since the root subalgebras generate \(\mathfrak g\), both actions are trivial. The proposition concerns complex semisimple Lie algebras and the listed dimension triples.

\section{Spinorial restrictions and compatible subgroups}\label{sec:spin}

\subsection{The two spin identities}\label{subsec:D1}

For a rank-\(r\) spin group, use torus variables in which the vector weights are \(\pm2e_i\) and the half-spin weights are the vectors \((\epsilon_1,\ldots,\epsilon_r)\) with \(\epsilon_i\in\{\pm1\}\) and a fixed parity of the number of minus signs. If \(s_\pm(x)\) is the corresponding sum of Laurent monomials, the exterior-square character is

\[
\frac{s_\pm(x)^2-s_\pm(x^2)}2.
\]

Expanding finite Laurent polynomials at \(r=5,6\) gives respectively

\[
\Lambda^2S_{10}^{\pm}\cong\Lambda^3V_{10},\qquad
\Lambda^2S_{12}^{\pm}\cong\Lambda^4V_{12}\oplus\mathbf1.
\]

The program \path{verify_spin_arithmetic.py} compares the Laurent coefficients on a maximal torus using the weight sets above. Maximal-torus characters determine these compact-group representations. Baez credits Schwahn with the explanation that motivated this article \citep{Baez2026Coincidences}.

\subsection{The quadratic model and its spin lift}\label{subsec:D2}

\paragraph{The covering theorem and the sign of a full turn.}
We use the classical theorem that, for \(d\geq3\), \(\operatorname{Spin}(d)\) is the simply connected double cover of \(\mathrm{SO}(d)\), with kernel \(\{1,-1\}\); see Gallier \citep[Theorem 1.11 and Section 1.8]{Gallier2014Spin}. In particular \(\pi_1(\mathrm{SO}(d))\cong\mathbb Z/2\), generated by a full rotation in one coordinate plane. This topological fact is the lifting criterion used below.

To describe the lifted rotation in a coordinate plane, use the real Euclidean Clifford algebra with \(e_i^2=-1\), \(e_ie_j=-e_je_i\) for \(i\ne j\), and put \(B=e_1e_2\), so \(B^2=-1\). The path
\[
s(\theta)=\exp(\theta B/2)
=\cos(\theta/2)+B\sin(\theta/2)
\]
acts on vectors by conjugation. Direct multiplication gives
\[
s(\theta)e_1s(\theta)^{-1}
   =\cos\theta\,e_1+\sin\theta\,e_2,\qquad
s(\theta)e_2s(\theta)^{-1}
   =-\sin\theta\,e_1+\cos\theta\,e_2.
\]
It fixes the other coordinate vectors. At \(2\pi\) it is \(-1\), and at \(4\pi\) it is \(1\). Bivectors belonging to disjoint coordinate planes commute; hence a loop with integer winding numbers in the coordinate planes \(r_1,\ldots,r_s\) has endpoint \((-1)^{r_1+\cdots+r_s}\).

A continuous homomorphism from \(\mathrm{SO}(4)\) to \(\mathrm{SO}(10)\) lifts through this double cover precisely when its induced map on the fundamental group is zero. Write \(\widetilde\rho\) for the lift taking the identity to the identity. On the connected space \(\mathrm{SO}(4)\times\mathrm{SO}(4)\), the maps \((g,h)\mapsto\widetilde\rho(gh)\) and \((g,h)\mapsto\widetilde\rho(g)\widetilde\rho(h)\) lift the same map and agree at the identity; uniqueness of lifts makes them equal, proving that \(\widetilde\rho\) is a homomorphism. The angle computation below verifies this criterion for the symmetric-square action.

Let \(A_{\mathbb R}=\{x\in\mathbb R^5:\sum_i x_i=0\}\) and let \(A\) be its complexification. Restrict the ten off-diagonal monomials \(x_i x_j\), \(i<j\), to this hyperplane. They span its quadratic forms because \(x_i^2=-\sum_{j\ne i}x_i x_j\). If an off-diagonal quadratic vanishes there, it is divisible by the linear equation of the hyperplane, hence has the form \((\sum_i x_i)(\sum_i a_i x_i)\). Its diagonal coefficients are \(a_i\), and consequently every coefficient is zero. This identifies the edge permutation representation with \(\operatorname{Sym}^2 A^*\), and the invariant metric identifies \(A^*\) with \(A\).

The symmetric-square representation of \(\mathrm{SO}(4)\) has plane angles

\[
2x,\quad 2y,\quad x+y,\quad x-y,\quad 0
\]

on a maximal torus. A generator of its fundamental group rotates one source plane through \(2\pi\), so the sum of the winding numbers in the target rotation planes is \(2+1+1=0\) modulo two. The homomorphism into \(\mathrm{SO}(10)\) therefore lifts to \(\operatorname{Spin}(10)\). The metric is an invariant vector in the symmetric square, so the connected lifted group is contained in its stabilizer \(\operatorname{Spin}(9)\).

The kernel of the symmetric-square homomorphism is exactly \(\{I,-I\}\): if every \(v^2\) is fixed, each \(v\) is carried to \(\pm v\), and linearity applied to sums forces the same sign for all nonzero vectors. To compute the lift at \(-I\), rotate both source planes from zero to \(\pi\). Three target planes make full turns, so the path lifts to the nonidentity central element \(-1\). The lifted homomorphism is therefore injective, although its orthogonal representation has kernel \(\{I,-I\}\).

\subsection{The half-spin representation as a kernel}\label{subsec:D3}

On an oriented Euclidean four-space define

\[
w(a\otimes(b\wedge d))=a\wedge b\wedge d,\qquad
c(a\otimes(b\wedge d))=\langle a,b\rangle d-\langle a,d\rangle b.
\]

With the usual orthonormal exterior bases, \(ww^*=cc^*=3I\) and \(wc^*=cw^*=0\). Therefore \((w,c)\) is onto \(\Lambda^3A\oplus A\), and \(\mathcal T(A)=\ker(w,c)\) has dimension \(24-8=16\). Since \(\Lambda^3A\cong A\) for the orientation-preserving group,

\[
\chi_{\mathcal T(A)}=\chi_A(\chi_{\Lambda^2A}-2).
\]

The half-spin characters for the lifted action above are

\[
\begin{aligned}
\chi_{S_{10}^{\pm}}(x,y)
&=16\cos x\cos y\cos\frac{x+y}{2}\cos\frac{x-y}{2}\\
&=8\cos x\cos y(\cos x+\cos y).
\end{aligned}
\]

The zero fifth plane angle makes the two half-spin characters equal on this subgroup. On the source,

\[
\chi_A=2\cos x+2\cos y,\qquad
\chi_{\Lambda^2A}=2+4\cos x\cos y.
\]

Thus the characters coincide on the entire maximal torus, proving \(S_{10}^{\pm}|_{\mathrm{SO}(4)}\cong\mathcal T(A)\). On the five \(A_5\) classes, the kernel has character \((16,0,-2,1,1)\), which identifies it with \(S=2B\oplus C\). The restricted vector representation is \(E=\mathbf1\oplus A\oplus B\).

The adjoint identities above prove surjectivity. The full-torus weight calculation is reproduced by \path{verify_spin_arithmetic.py}.

\paragraph{Joint surjectivity of wedge multiplication and contraction.}
On the orthonormal real tensor bases, for \(a<b<c\),
\[
\begin{aligned}
w^*(e_a\wedge e_b\wedge e_c)
&=e_a\otimes(e_b\wedge e_c)-e_b\otimes(e_a\wedge e_c)
   +e_c\otimes(e_a\wedge e_b),\\
c^*(v)&=\sum_{i=1}^4e_i\otimes(e_i\wedge v).
\end{aligned}
\]
Here star is the transpose for those real inner products, extended as the Hermitian adjoint after complexification; the contraction map itself is the complex-linear extension of the real bilinear formula. These expressions show \(ww^*=cc^*=3I\) and \(wc^*=cw^*=0\). In particular a preimage of \((\eta,v)\in\Lambda^3A\oplus A\) is \((w^*\eta+c^*v)/3\). This formula gives an explicit right inverse to \((w,c)\), with normalization factor \(1/3\), and proves its surjectivity. The elementary checks reconstruct the two matrices and verify these identities over the integers.

\subsection{The principal rotation subgroup and the compatible intersection}\label{subsec:D4}

Let \(H_j\) be the complexification of the real degree-\(j\) spherical harmonics, with torus weights \(-j,-j+1,\ldots,j\). The real nine-dimensional \(H_4\) gives a faithful representation of \(\mathrm{SO}(3)\) into \(\mathrm{SO}(9)\). Its winding parity is \(1+2+3+4=0\) modulo two, so it lifts. The lift is faithful because its composition with the covering map \(\operatorname{Spin}(9)\to\mathrm{SO}(9)\) is faithful.

The sixteen spinor weights are \(\frac12(\pm1\pm2\pm3\pm4)\); their multiplicity distribution is the union of the weights of \(H_5\) and \(H_2\). Hence

\[
V_{10}|_{\mathrm{SO}(3)}=H_0\oplus H_4,\qquad
S_{10}^{\pm}|_{\mathrm{SO}(3)}=H_2\oplus H_5.
\]

The harmonic character formula \(\chi_{H_j}(\theta)=1+2\sum_{m=1}^j\cos(m\theta)\) gives \(\chi_{H_4}=(9,1,0,-1,-1)\) on the icosahedral classes, equal to the character of \(A\oplus B\). Thus both the tensor and principal subgroup constructions restrict on their nine-dimensional vector spaces to the same real orthogonal \(A_5\)-module. Choose an orthogonal identification. If it reverses orientation, multiplying it by \(-I\) in dimension nine corrects the orientation without changing equivariance. Lift the identification into \(\operatorname{Spin}(9)\). The two lifts on \(A_5\) then agree, since their difference would be a character \(A_5\to\{1,-1\}\), and no such nontrivial character exists.

For these compatible images \(K_{\mathrm t}\) and \(K_{\mathrm p}\), the Lie algebra of the intersection is an \(A_5\)-submodule of the irreducible three-dimensional \(\mathfrak{so}(3)\). A positive-dimensional intersection would contain the entire principal group. At the Lie algebra level the tensor group has two \(\mathfrak{so}(3)\) factors and acts on \(\mathbf3\otimes\mathbf3\). A nonzero homomorphism from \(\mathfrak{so}(3)\) to a factor is an isomorphism. Consequently a connected rotation subgroup acts either through one factor, giving three copies of \(\mathbf3\), or through both, giving \(\mathbf1\oplus\mathbf3\oplus\mathbf5\). Neither is the irreducible \(H_4\), so the intersection is finite. The finite rotation groups are cyclic, dihedral, tetrahedral, octahedral, or icosahedral; one containing \(A_5\) is \(A_5\). Therefore \(K_{\mathrm t}\cap K_{\mathrm p}=A_5\).

The conclusion \(K_{\mathrm t}\cap K_{\mathrm p}=A_5\) holds for the compatible subgroup images chosen above.

To test whether the \(210\) identity extends to the full spin group, take an element \(g\) whose vector eigenvalues are \(-1,-1,1^8\). On \(\mathbf1\oplus2V_{10}\), the element \(g\) has trace thirteen and \(g^2\) has trace twenty-one, so \(g\) has trace \(74\) on \(\Lambda^2(\mathbf1\oplus2V_{10})\). The coefficient of \(t^4\) in \((1-t)^2(1+t)^8\) is \(-14\). These full spin-group representations are not isomorphic even though their restrictions to the marked \(A_5\) are isomorphic.

\subsection{Polynomial harmonics: dimension, weights, and irreducibility}\label{subsec:D5}
Let \(\mathcal P_j\) be the homogeneous degree-\(j\) polynomials in three real variables with complex coefficients, and set \(\mathcal P_j=0\) for \(j<0\). Put \(R=x^2+y^2+z^2\) and \(H_j=\ker(\Delta:\mathcal P_j\to\mathcal P_{j-2})\). The function action is \((gp)(v)=p(g^{-1}v)\). The rotation action fixes \(R\) and commutes with \(\Delta\).

Give monomials the positive-definite Hermitian inner product
\[
\langle x^ay^bz^c,x^{a'}y^{b'}z^{c'}\rangle
=\begin{cases}a!b!c!&(a,b,c)=(a',b',c'),\\0&\text{otherwise}.\end{cases}
\]
Differentiation with respect to a variable is adjoint to multiplication by that variable. Hence the adjoint of \(\Delta:\mathcal P_j\to\mathcal P_{j-2}\) is multiplication by \(R\). Since multiplication by the nonzero polynomial \(R\) is injective, \(\Delta\) is onto. Finite-dimensional orthogonal decomposition gives
\[
\mathcal P_j=H_j\oplus R\mathcal P_{j-2},\qquad
\dim H_j=\binom{j+2}{2}-\binom j2=2j+1
\quad(j\geq2).
\]
Both summands are invariant under real rotations. \(H_0\) consists of the constants and \(H_1\) of the linear polynomials, so their dimensions are \(1\) and \(3\). Restriction to the sphere is injective on homogeneous polynomials, since scaling a nonzero point to the sphere recovers its value.

For a rotation about the \(z\)-axis, put \(u=x-iy\), \(v=x+iy\). Under the chosen function action, \(u,v,z\) have weights \(1,-1,0\). A monomial \(u^av^bz^c\), \(a+b+c=j\), has weight \(a-b\). For \(|m|\leq j\), the multiplicity of weight \(m\) in \(\mathcal P_j\) is
\[
\left\lfloor\frac{j-|m|}{2}\right\rfloor+1;
\]
it is zero otherwise. Multiplication by \(R=uv+z^2\) preserves weight. Subtracting the multiplicities for \(\mathcal P_{j-2}\) leaves exactly one of each weight \(-j,-j+1,\ldots,j\) in \(H_j\).

For completeness, irreducibility can also be seen at the polynomial level. In these coordinates \(\Delta=4\partial_u\partial_v+\partial_z^2\). The operators
\[
H=u\partial_u-v\partial_v,\qquad
E=u\partial_z-2z\partial_v,\qquad
F=2z\partial_u-v\partial_z
\]
are complex linear combinations of infinitesimal real rotations. They preserve degree, commute with \(\Delta\), and satisfy
\[
[H,E]=E,\qquad[H,F]=-F,\qquad[E,F]=2H.
\]
The polynomial \(u^j\) belongs to \(H_j\), is killed by \(E\), and has \(H\)-eigenvalue \(j\). Commuting \(E\) past powers of \(F\) yields
\[
E F^r u^j=r(2j-r+1)F^{r-1}u^j.
\]
Thus \(F^r u^j\ne0\) for \(0\leq r\leq2j\), by induction, and these vectors have the distinct weights \(j-r\). They form a basis by the dimension computation. A nonzero rotation-invariant subspace is invariant under \(H,E,F\). Polynomials in the diagonal operator \(H\) project onto its weight spaces, so the invariant subspace contains a nonzero weight vector; repeated application of \(E\) reaches \(u^j\), and then powers of \(F\) recover every basis vector. This proves irreducibility. The explicit quadratic basis gives a concrete model for these operators, while the argument above proves the statement uniformly in \(j\).

\paragraph{Why these are all the irreducible rotation representations.}
A finite-dimensional continuous representation of the compact rotation group has an invariant positive-definite Hermitian inner product by averaging over the group, and therefore splits into irreducibles. In an irreducible representation the circle of rotations about one axis has integral weights, since a full turn is the identity. Choose a vector \(v\) of maximal weight \(j\). The complexified infinitesimal operators \(H,E,F\) have the commutators displayed above, so \(Ev=0\) and
\[
E F^r v=r(2j-r+1)F^{r-1}v.
\]
Let \(r>0\) be the first index with \(F^rv=0\), which exists because weights are bounded below. The formula and \(F^{r-1}v\ne0\) give \(r=2j+1\); in particular \(j\) is a nonnegative integer. The span of \(v,Fv,\ldots,F^{2j}v\) is invariant under the three infinitesimal operators, hence under the connected rotation group, and irreducibility implies that this nonzero invariant subspace is the whole space. These operators act with the same coefficients as on the harmonic basis constructed above, giving an isomorphism with \(H_j\). Thus every irreducible representation is one of the \(H_j\), as used in Section F.4.

Finally, the diagonal action on \(\mathbb C^3\otimes\mathbb C^3\) has the scalar, skew-symmetric, and symmetric trace-free matrix summands. The cross-product map identifies the skew-symmetric part with \(H_1\), and \(M\mapsto x^{\mathsf T}Mx\) identifies the symmetric trace-free part with \(H_2\), since its Laplacian is \(2\operatorname{tr}M\). This proves the \(H_0\oplus H_1\oplus H_2\) decomposition used in Section D.4. The two rotation factors of \(\mathrm{SO}(4)\) can also be obtained from its action on self-dual and anti-self-dual two-forms. Their nine-dimensional tensor product is the trace-free symmetric-square representation. Thus the decompositions are representation-theoretic, not merely numerical.

\section{Arithmetic of the repetitions}\label{sec:arithmetic}

\subsection{The elliptic equation and its point additions}\label{subsec:E1}

Write the equation as \(\binom m2=\binom n3\), with indices \(m\) and \(n\). Setting \(x=3(n-1)\), \(y=9m-5\) gives the exact polynomial identity

\[
y^2+y-x^3+9x-20
=162\left(\binom m2-\binom n3\right).
\]

The generalized Weierstrass model \(y^2+y=x^3-9x+20\) has discriminant \(-3^6\cdot179\ne0\). To recover integral indices from an integral point, one requires \(x\equiv0\pmod3\) and \(y\equiv4\pmod9\), together with the positivity and half-triangle conditions on \((m,n)\) used in the article.

The negative of \((x,y)\) in the group law is \(-(x,y)=(x,-y-1)\). For points \((x,y),(x',y')\) that are not negatives of one another, the chord or tangent slope is

\[
\lambda=\frac{y'-y}{x'-x},\qquad
\lambda=\frac{3x^2-9}{2y+1}\quad\text{when the points coincide},
\]

and their sum is

\[
x''=\lambda^2-x-x',\qquad
y''=-1-y+\lambda(x-x'').
\]

With \(P_0=(0,4)\) and \(Q_0=(3,4)\), the exact additions give

\[
2P_0=(1,-4),\qquad
2Q_0-P_0=(27,139),\qquad
3P_0=(63,499),\qquad
-(P_0+3Q_0)=(105,1075).
\]

The tangent at \(P_0\) is \(y=4-x\), and its third intersection with the cubic, counted with multiplicity, is \((1,3)\). The chord through \(P_0\) and \(2P_0\) is \(y=4-8x\); substitution into the curve equation gives \(0=-x(x-1)(x-63)\), so its third intersection with the cubic is \((63,-500)\). Negating this point in the group law gives the displayed value of \(3P_0\). The three listed points recover \((n,m)=(10,16),(22,56),(36,120)\) and the values \(120,1540,7140\).

De Weger \citep[Section 3, p. 377]{deWeger1997} gives the classical elliptic-curve interpretation of the binomial equation. The coordinates and point combinations above are our calculation in this Weierstrass model.

Avanesov's classification of positive common triangular and tetrahedral values is
\[
1,\quad10,\quad120,\quad1540,\quad7140.
\]
We use this as a classical theorem \citep{Avanesov1967}, separate from the explicit group-law calculations above.

\subsection{The two fifth-order defect identities}\label{subsec:E2}

Let \(t\) be positive and odd and \(r=(t+1)/2\). For \(n=15t^2\), set

\[
m=\frac{t(n^2-5n+5)+1}{2},\quad
u=\frac t2(n-1)(n-4),\quad
v=\frac t2(n-2)(n-3).
\]

Then \(v-u=t\), \(u=m-r\), \(v=m+r-1\), and

\[
uv=2\binom n5=m(m-1)-r(r-1).
\]

For \(n=15t^2+4\), set

\[
m=\frac{t(n^2-3n+1)+1}{2},\quad
u=\frac t2n(n-3),\quad
v=\frac t2(n-1)(n-2).
\]

The same relations hold. All the displayed quantities are integers in the stated range. Dividing the last equality by two proves in both cases

\[
\binom m2-\binom n5=\binom r2.
\]

At \(t=1\) this yields \((n,m)=(15,78)\) and \((19,153)\). For odd \(t>1\) the defect is positive. These are the published families \citep[Section 3.2, (3)--(4)]{BlokhuisBrouwerdeWeger2017} with its parameter replaced by \(r=(t+1)/2\). In that paper the corresponding quintic expressions for the index \(m\) in \(\binom m2\) are

\[
3600r^5-9000r^4+8700r^3-4050r^2+905r-77
\]

and

\[
3600r^5-9000r^4+9300r^3-4950r^2+1355r-152.
\]

Substituting \(t=2r-1\) recovers the two polynomial families in the cited paper. The present formulas are simply a reparametrization that makes their factorization convenient.

A nearby spinorial identity is \(\Lambda^2S_{12}^{\pm}\cong\Lambda^4V_{12}\oplus\mathbf1\). It gives \(496=495+1\), as noted in Baez's motivating account \citep{Baez2026Coincidences}. The verification script compares both chiralities on the full rank-six torus.

\subsection{The normalized Pell–Fibonacci statement}\label{subsec:E3}

Require \(k\ge2\) and \(n\ge2k+3\), so that both entries of \(\binom nk=\binom{n-1}{k+1}\) lie in the normalized nontrivial half-triangle. Cancelling factorials gives

\[
n(k+1)=(n-k)(n-k-1),\qquad
n^2-(3k+2)n+k^2+k=0.
\]

For integers \(n,k\) in the stated range, the equation holds if and only if there is a positive integer \(d\) satisfying

\[
d^2=5k^2+8k+4,\quad d>0,\quad
n=\frac{3k+2+d}{2}.
\]

Since \(d^2-(k+2)^2=4k(k+1)>0\), the other root is smaller than \(k\). Parity in the discriminant equation ensures that the displayed positive root is an integer when \(d\) is. Completing the square gives

\[
(5k+4)^2-5d^2=-4.
\]

Lind's theorem gives all positive solutions in terms of the Fibonacci and Lucas sequences \citep[Theorem 7 and Section 5]{Lind1968}. The norm identity is \(L_j^2-5F_j^2=4(-1)^j\). Thus the negative norm requires odd \(j\), and the residue pattern of the Lucas numbers modulo five, \(2,1,3,4\) with period four, gives \(L_j\equiv4\pmod5\) exactly when \(j\equiv3\pmod4\). Write \(j=4i+3\). The case \(i=0\) is a boundary case; the normalized family has \(i\ge1\) and

\[
n_i=F_{2i+2}F_{2i+3},\qquad
k_i=F_{2i}F_{2i+3}.
\]

Lind solves the quadratic in \(k\), obtaining \((5n+1)^2-5v^2=-4\); his parametrization agrees with ours after the index change \(s=i+1\) and the displayed change of variables. In our parametrization, the value of \(n\) is given by the positive-root formula above.

The first pairs are \((15,5),(104,39),(714,272),(4895,1869)\). The script \path{verify_spin_arithmetic.py} checks the initial members. Completeness comes from Lind's solution of the Pell equation.

\paragraph{Sequence identities and recovery of both indices.}
Let \(\alpha=(1+\sqrt5)/2\), \(\beta=(1-\sqrt5)/2\); then \(\alpha\beta=-1\), \(\alpha-\beta=\sqrt5\). Solving the elementary recurrence gives
\[
F_j=\frac{\alpha^j-\beta^j}{\sqrt5},\qquad L_j=\alpha^j+\beta^j.
\]
For integers \(b\geq a\geq0\), multiplication and \(\alpha\beta=-1\) give
\[
5F_aF_b=L_{a+b}-(-1)^aL_{b-a},\qquad
L_aF_b=F_{a+b}+(-1)^aF_{b-a}.
\]
Set \(a=2i,b=2i+3\). The first identity gives
\(5F_{2i}F_{2i+3}=L_{4i+3}-L_3=L_{4i+3}-4\).
For \(a\geq2\), the recurrence also gives \(2F_{a+2}-3F_a=L_a\), so the second identity yields
\[
2F_{2i+2}F_{2i+3}-3F_{2i}F_{2i+3}
=F_{4i+3}+F_3=F_{4i+3}+2.
\]
Together with Lind's description of the positive Pell solutions, these identities yield the displayed family from the positive-root formula for \(n\) and the congruence formula for \(k\).

The Lucas residues modulo five are \(2,1,3,4\) periodically: after four steps the ordered pair of consecutive residues returns to \((2,1)\), which determines every later residue by the recurrence. The congruence \(L_j\equiv4\pmod5\) therefore selects \(j\equiv3\pmod4\). At \(j=3\), \(k=0\) is a boundary case; at \(j=7\), \(L_7=29,F_7=13\) give \(k=5,n=15\). For positive \(k\), \(d^2>(k+2)^2\) implies \(n>2k+2\), so the positive integral root satisfies the half-triangle condition.

\section{Klein residues and the factored intertwiner}\label{sec:residues}

\subsection{The divisor and its linearizations}\label{subsec:F1}

Let \(G=A_5\), identified with the rotations of the icosahedron, and let \(\widetilde G\subset\mathrm{SU}(2)\) be its inverse image under \(\mathrm{SU}(2)\to\mathrm{SO}(3)\). We use the even binary-form representations

\[
H_j=\operatorname{Sym}^{2j}\mathbb C^2
    \cong H^0(\mathbb P^1,\mathcal O(2j)),\qquad j\ge0.
\]

An invariant bilinear pairing identifies the representation and its dual where needed. Even degree ensures that the central involution acts trivially. The coordinate and binary-action conventions are fixed below. Klein's invariant form is

\[
F(X,Y)=X^{11}Y+11X^6Y^6-XY^{11}.
\]

The group fixes \(F\) itself, and its twelve distinct zeros form the reduced divisor \(D\), each point having multiplicity one. For the classical form and transvectants, see Lombardo--Sanders \citep[Sections 3 and 8.1]{LombardoSanders2010}; our derivative ordering in Section~\ref{subsec:F5} differs from theirs by the factor $(-1)^r$ for the $r$th transvectant. In the affine chart \(z=X/Y\), put

\[
f(z)=z^{11}+11z^6-z.
\]

The polynomial is squarefree, and the homogeneous form has one further simple zero at infinity. The spaces

\[
I_j=H^0(D,\mathcal O(2j)|_D)
\]

have dimension twelve, with the action on the fibers inherited from the line bundle. Each \(I_j\) is induced from the character by which the vertex stabilizer \(C_5\) acts on the fiber; that character depends on \(j\bmod5\).

\subsection{Division, differentiation, and low-degree residues}\label{subsec:F2}

Tensoring the divisor sequence by \(\mathcal O(2j)\) gives

\[
0\longrightarrow\mathcal O(2j-12)
\xrightarrow{\ F\ }\mathcal O(2j)
\longrightarrow\mathcal O(2j)|_D\longrightarrow0.
\]

For \(j\ge6\), the vanishing of \(H^1(\mathcal O(2j-12))\) yields

\[
0\longrightarrow H_{j-6}\xrightarrow{F}H_j
\longrightarrow I_j\longrightarrow0.
\]

The kernel of evaluation has dimension \(2j-11\), the target \(I_j\) has dimension twelve, and \(\dim H_j=2j+1\). All maps are equivariant because \(F\) is fixed. We choose the usual invariant positive-definite Hermitian inner products on the harmonic spaces and use orthogonal complements to split these finite-dimensional sequences.

The first derivative of a line-bundle section on its zero set is intrinsic: if the local defining function is multiplied by a nowhere-vanishing transition function, its derivative on the zero set changes by that transition alone. Consequently

\[
\delta=dF|_D\in H^0(D,K_{\mathbb P^1}\otimes\mathcal O(12)|_D)
=H^0(D,\mathcal O(10)|_D).
\]

It is invariant and nowhere zero, so multiplication gives \(I_j\cong I_{j+5}\). This trivializes \(\mathcal O(10)\) along the divisor \(D\). In the compatible fiber coordinates used for the rational matrices, multiplication is represented by \(f'(z)\) at the finite roots and by \(-1\) at infinity. The latter sign incorporates the differential transition between \(z\) and \(1/z\).

For \(0\le r\le5\), the same divisor sequence and equivariant Serre duality instead give

\[
0\longrightarrow H_r\longrightarrow I_r
\xrightarrow{\beta_r}H_{5-r}^{*}\longrightarrow0.
\]

The map \(\beta_r\) is given by a residue pairing. For \(a\in I_r\) and \(b\in H_{5-r}\), the product \(a\,b|_D\) is a section of \(\mathcal O(10)|_D\); division by \(F\) and the differential identification give the following formula in affine coordinates:

\[
\beta_r(a)(b)
=\sum_{f(\alpha)=0}\frac{a(\alpha)b(\alpha)}{f'(\alpha)}
-a_\infty b_\infty.
\]

For any polynomial \(p\), the finite-root sum \(\sum p(\alpha)/f'(\alpha)\) is the coefficient of \(z^{10}\) in the remainder of \(p\) modulo \(f\). This follows by Lagrange interpolation for the monic polynomial of degree eleven. If \(a\) is the restriction of a form in \(H_r\), the coefficient of \(z^{10}\) in the product is \(a_\infty b_\infty\), so the finite-root sum cancels the contribution at infinity. Thus \(\beta_r\) annihilates the restriction image. Its rank is \(11-2r\), and the image of \(H_r\) has dimension \(2r+1\).

At \(r=0\), this is the pairing of functions on twelve points with sum-zero residue vectors. After choosing an invariant splitting and identifying \(H_5^*\) with \(H_5\) by its invariant bilinear pairing, it gives \(I_0=Q\cong H_0\oplus H_5\). This realization of \(Q\) uses the residue sequence; the realization of \(P\) below uses evaluation of quadratic forms.

The model $P\cong H_0\oplus H_2$ is obtained by evaluating real quadratic forms on six unit representatives of the antipodal axes. The associated rank-one symmetric tensors have Gram matrix

\[
\frac45 I_6+\frac15\mathbf 1\mathbf 1^{\mathsf T},
\]

with eigenvalues \(2\) and \(4/5\) (the latter five times), so evaluation is invertible. The trace and trace-free parts of the quadratic forms are \(H_0\) and \(H_2\).

\subsection{The positive reduction and the factored map}\label{subsec:F3}

For \(j\ge6\), set \(r=j\bmod5\) with \(0\le r\le4\). Split the high-degree sequence, multiply its quotient by \(\delta^{-(j-r)/5}\), and split the low-degree sequence. These choices define an equivariant isomorphism

\[
T_j:H_j|_G\longrightarrow
H_{j-6}|_G\oplus H_r|_G\oplus H_{5-r}|_G.
\]

The remaining harmonic summand has degree \(j-6\), and the other summands are among \(H_0,\ldots,H_5\). After restriction to \(A_5\), these examples give \(H_6\cong H_0\oplus H_1\oplus H_4\) and \(H_8\cong H_3\oplus2H_2\).

The two rotation-group spaces

\[
\widetilde L=\Lambda^2\operatorname{Sym}^2(H_0\oplus H_5),\qquad
\widetilde R=\Lambda^5\Lambda^2(H_0\oplus H_2)
\]

have the following harmonic multiplicities:

{\small
\begin{longtable}[]{@{}rrr@{}}
\toprule\noalign{}
\(j\) & \(\widetilde L\) & \(\widetilde R\) \\
\midrule\noalign{}
\endhead
\bottomrule\noalign{}
\endlastfoot
0 & 1 & 14 \\
1 & 8 & 48 \\
2 & 8 & 60 \\
3 & 14 & 72 \\
4 & 13 & 61 \\
5 & 19 & 53 \\
6 & 16 & 32 \\
7 & 18 & 20 \\
8 & 15 & 8 \\
9 & 16 & 3 \\
10 & 13 & 0 \\
11 & 12 & 0 \\
12 & 8 & 0 \\
13 & 8 & 0 \\
14 & 5 & 0 \\
15 & 5 & 0 \\
16 & 2 & 0 \\
17 & 2 & 0 \\
18 & 1 & 0 \\
19 & 1 & 0 \\
\end{longtable}
}

These multiplicities can be recovered from torus weights by subtracting the multiplicity of the weight \(j+1\) from that of the weight \(j\). Applying the positive reduction to both columns produces

\[
\mathcal R=46H_0\oplus103H_1\oplus96H_2
\oplus103H_3\oplus96H_4\oplus53H_5.
\]

Choose Clebsch--Gordan decompositions and fix an order on the repeated summands produced by each chosen reduction. Together with the invariant splittings, these choices define

\[
\Phi_L:\Lambda^2\operatorname{Sym}^2Q\longrightarrow\mathcal R|_G,
\qquad
\Phi_R:\Lambda^5\Lambda^2P\longrightarrow\mathcal R|_G.
\]

Then $I_F=\Phi_R^{-1}\Phi_L$ is a noncanonical equivariant isomorphism. The reductions in this construction have matrices with at most \(39\) rows and \(39\) columns. The reduction matrices are in \texttt{verification/prior\_inputs/}; \path{verify_geometry.py} checks their exact ranks and two-generator equivariance. The factorization specifies the intertwiner using these smaller matrices; the chosen splittings and multiplicity-space identifications are part of its definition.

For \(210\), the same procedure applied to \(\Lambda^2(3H_0\oplus2H_4)\) and \(\Lambda^4(H_0\oplus H_4)\) yields

\[
5H_0\oplus7H_1\oplus6H_2\oplus7H_3\oplus8H_4\oplus3H_5.
\]

For the \(3003\) identity, Section B gives a proof using signed orbits and shows that the two unsigned bases are not isomorphic as \(G\)-sets.

\subsection{The kernel of harmonic restriction}\label{subsec:F4}

Let \(t=[H_1]\). For \(j\ge1\), the tensor rule gives

\[
[H_{j+1}]=(t-1)[H_j]-[H_{j-1}],\qquad [H_0]=1,\quad[H_1]=t,
\]

and hence \(R(\mathrm{SO}(3))=\mathbb Z[t]\). The five rotation-character values on \(G\) are \(3,-1,0,\varphi,1-\varphi\). Their monic vanishing polynomial is

\[
p(t)=t(t+1)(t-3)(t^2-t-1).
\]

A representation-ring element restricts to zero precisely when its character vanishes on all five classes. Dividing its polynomial by the monic \(p\) over \(\mathbb Z\) leaves a remainder of degree at most four, which cannot vanish at all five distinct roots unless it is zero. Thus the kernel of restriction is the principal ideal \((p)\). Equivalently,

\[
[H_5]+[H_4]-[H_3]-2[H_2]-[H_1]=p(t).
\]

Hence the kernel of harmonic restriction on these six generators is generated by the displayed relation.

\subsection{Matrix formulas for the chosen splittings}\label{subsec:F5}

The matrices below implement the chosen splittings and identifications explicitly. In the affine monomial basis ordered by increasing exponent, the standard binary-form Hermitian metric on \(H_j\) is diagonal with entries \(1/\binom{2j}{i}\). Write \(E_j\) for restriction to the eleven finite-root coordinates modulo \(f\) together with the coordinate at infinity. For a surjective restriction in degree \(s\), the quotient metric is

\[
G_s=(E_sD_s^{-1}E_s^*)^{-1}.
\]

For \(r=0,\ldots,4\), use \(s=r+10\) and transport the quotient metric to \(I_r\) by multiplication by \(\delta^2\). Orthogonal projection gives the left inverse of $H_r\to I_r$; combine it with the residue map \(\beta_r\) and the invariant bilinear self-duality on \(H_{5-r}\). The splitting in high degree is the orthogonal projection along \(FH_{j-6}\). This specifies \(T_j:H_j\to H_{j-6}\oplus H_r\oplus H_{5-r}\) with summands in the displayed order. The archive contains these matrices.

The fourteen matrices for \(6\leq j\leq19\) have full rank over \(\mathbb Q\). Their equivariance can be checked on a diagonal fifth-root generator and the second binary generator

\[
\frac{i}{\sqrt{1+\varphi^2}}
\begin{pmatrix}1&\varphi\\ \varphi&-1\end{pmatrix},
\qquad\varphi=\frac{1+\sqrt5}{2}.
\]

The latter acts on the even binary forms over \(\mathbb Q(\sqrt5)\); both coefficient components of every intertwining equation are checked exactly. In the chosen fiber coordinates the determinant of multiplication by \(\delta\) is \(5^{25}\), including the infinity factor. These finite tests supplement the intrinsic equivariance and splitting arguments above.

The unnormalized transvectant convention used below is

\[
(f,g)_r=\sum_{a=0}^{r}(-1)^a\binom ra
\frac{\partial^rf}{\partial X^{r-a}\partial Y^a}
\frac{\partial^rg}{\partial X^a\partial Y^{r-a}}.
\]

The transvectant is equivariant for the simultaneous action on both arguments. The map

\[
H_5\oplus H_4\longrightarrow H_1\oplus H_3\oplus H_2\oplus H_2,
\quad(u,v)\longmapsto
\bigl((F,u)_{10},(F,u)_8+(F,v)_7,(F,u)_9,(F,v)_8\bigr)
\]

is invertible: its determinant has absolute value \(2^{262}3^{125}5^{75}7^{32}11^{20}\). Reversing the monomial bases changes the sign; the certificate records ascending powers. In particular \(B\cong H_2\), \(A\cong\ker((F,-)_8:H_4\to H_2)\) and \(C\cong\ker((F,-)_9:H_5\to H_2)\). The ranks of the latter maps are five. These kernel descriptions identify \(A,B,C\) in the binary-form model.

\subsection{Polynomial models, local residues, and the elementary exactness proof}\label{subsec:F6}

Throughout this subsection we use the maps and normalizations fixed in Sections~\ref{subsec:F1}--\ref{subsec:F5}.

For evaluation at \(0,1,-1\) on polynomials of degree at most three, the matrix in the basis \((1,z,z^2,z^3)\) is
\[
E=\begin{pmatrix}1&0&0&0\\1&1&1&1\\1&-1&1&-1\end{pmatrix}.
\]
It has rank three and kernel spanned by \(z^3-z\). The three interpolation polynomials \(1-z^2,z(z+1)/2,z(z-1)/2\) give an explicit right inverse to evaluation.

The binary model can also be checked directly. Put \(d=2j\), and let \(p(X,Y)=\sum_{i=0}^d p_iX^iY^{d-i}\), acted on by inverse substitution from \(\mathrm{SU}(2)\). On its diagonal subgroup the weights are \(j-i\). The operators \(X\partial_Y\) and \(Y\partial_X\) move between neighboring monomials with nonzero coefficients until an endpoint is reached. An invariant subspace decomposes into weight spaces, and these operators force a nonzero such subspace to contain every monomial. Thus this model is irreducible. The center acts by \((-1)^d=1\), so it descends to \(\mathrm{SO}(3)\), and the weights identify it with the \(H_j\) of Section~\ref{subsec:D5}.

For even \(d\), the pairing
\[
B_d(p,q)=\sum_{i=0}^d\frac{(-1)^i}{\binom di}\,p_iq_{d-i}
\]
is symmetric and nondegenerate. It is the pairing obtained by contracting the symmetric tensors with the determinant form on \(\mathbb C^2\), and hence is \(\mathrm{SL}_2\)-invariant. This also follows by differentiating under the two operators just displayed and under their diagonal commutator. It identifies the dual binary representation with itself; it is complex bilinear, not the positive-definite Hermitian inner product used to choose complements.

The affine polynomial factors as \(f(z)=z h(z^5)\), with \(h(u)=u^2+11u-1\). The two roots of \(h\) are \(\varphi^{-5}\) and \(-\varphi^5\), where \(\varphi=(1+\sqrt5)/2\). At a nonzero root, \(f'(z)=5u(2u+11)\ne0\); at zero, \(f'(0)=-1\). The expression for \(F\) in the chart at infinity is \(f_\infty(w)=w+11w^6-w^{11}\), with \(f_\infty'(0)=1\). Hence the complete homogeneous divisor is squarefree, including the point at infinity.

There is a direct interpolation proof of the high-degree surjection. For \(d\geq11\) (hence \(j\geq6\) when \(d=2j\)), take the polynomial of degree at most ten \(I(z)\) with prescribed values at the eleven finite roots, and add
\[
a_\infty z^{d-11}f(z).
\]
The new term vanishes at every finite root and has coefficient \(a_\infty\) at \(z^d\), which is the prescribed value at \([1:0]\). Homogenization gives the required degree-\(d\) section. A section whose twelve values vanish is divisible by all twelve distinct linear factors of \(F\), proving the asserted kernel. The resulting exact sequence is intrinsic even though this particular right inverse is coordinate-dependent; the invariant orthogonal splitting in Section~\ref{subsec:F5} is the one used in the supplied operators.

For clarity, the derivative sign can be read in both charts. For a degree-twelve section,
\[
f_\infty(w)=w^{12}f(1/w),\qquad
f_\infty'(w)=12w^{11}f(1/w)-w^{10}f'(1/w).
\]
Along a zero away from the two poles, the first term vanishes. Identifying \(K\otimes\mathcal O(12)\) with \(\mathcal O(10)\) by the finite expression \(p(z)\,dz\) uses the expression \(-p_\infty(w)\,dw\) in the infinity frame. Thus the fiber scalar of \(\delta\) is \(f'(\alpha_i)\) at a finite root and \(-f_\infty'(0)=-1\) at infinity, agreeing with the matrices in Section~\ref{subsec:F5}. The sign has no bearing on nonvanishing, but it is essential in the weighted residue pairing.

Interpolation proves the low-degree exact sequences as follows. For \(\deg h\leq10\), the Lagrange identity is
\[
h(z)=\sum_{i=1}^{11}\frac{h(\alpha_i)}{f'(\alpha_i)}
                           \frac{f(z)}{z-\alpha_i}.
\]
Both sides have degree at most ten and agree at the eleven roots. Taking coefficients of \(z^{10}\) gives
\[
[z^{10}]h(z)=\sum_{i=1}^{11}\frac{h(\alpha_i)}{f'(\alpha_i)}.
\]
For arbitrary \(h\), replace it on the right by its remainder modulo \(f\). For \(a=\operatorname{ev}_r p\) and \(b\in H_{5-r}\), the product \(p(z,1)b(z,1)\) has degree at most ten, and its coefficient of \(z^{10}\) is \(p(1,0)b(1,0)\); hence \(\beta_r\operatorname{ev}_r=0\). To show that \(\beta_r\) is onto, suppose \(b\) is annihilated by all functionals \(\beta_r(a)\). Taking \(a_\infty=0\) and varying one finite value at a time forces \(b(\alpha_i,1)=0\) for every \(i\). Since its degree is at most \(10-2r<11\), it vanishes identically. The transpose of \(\beta_r\) is therefore injective, and \(\beta_r\) is surjective. The rank is \(11-2r\). Evaluation on \(H_r\), of degree \(2r<12\), is injective, and the image of evaluation and \(\ker\beta_r\) both have dimension \(2r+1\). This proves exactness for every \(0\leq r\leq5\).

Equivariance follows from the intrinsic residue construction. The homogeneous differential \(Y\,dX-X\,dY\) is preserved under determinant-one substitutions; the quotient of the product of the two forms, whose degrees sum to ten, by \(F\), multiplied by this differential, has total scaling degree zero. The same interpretation applies to its local residue values when the first argument is prescribed only on \(D\). Thus the coordinate formula for \(\beta_r\) gives an invariant pairing. The general cohomological description of this exact sequence is Serre duality, but its special case here is established by the preceding elementary argument.

The six-axis evaluation has an equally explicit matrix proof. For unit representatives \(u_i\), evaluation of a symmetric matrix \(M\) is \(\operatorname{Tr}(M u_i u_i^{\mathsf t})\). The inner products of the six projectors are \((u_i\cdot u_j)^2\), equal to one on the diagonal and \(1/5\) off it. Hence their Gram matrix is \((4/5)I+(1/5)\boldsymbol1\boldsymbol1^{\mathsf t}\), with eigenvalues two and \(4/5\). They form a basis of the six-dimensional space of symmetric matrices, so evaluation is an equivariant isomorphism. These evaluation and residue maps give the geometric realizations of \(P\) and \(Q\).

Once the splittings are fixed, the component of the reduction on \(H_j\) sends \(p=Fq+s\) to \(q\) and the two components obtained from
\[
\delta^{-(j-r)/5}\operatorname{ev}_j(s)\in I_r,
\qquad r\equiv j\pmod5,\quad0\leq r\leq4.
\]
The projection onto the low-degree evaluation image and \(\beta_r\), followed by the bilinear self-duality, give those two components. Each constituent map is invertible on the indicated complementary pieces. At \(j=6\), the output is \(H_0\oplus H_1\oplus H_4\); at \(j=8\), \(H_3\oplus2H_2\); and the two-step reduction of \(H_{13}\) gives \(H_1\oplus2H_2\oplus2H_3\). They are restricted \(A_5\)-identities. Ordering the repeated copies gives the factored map of Section~\ref{subsec:F3}.

\subsection{Rotation-group models and their restrictions}\label{subsec:F7}

For the fixed icosahedral inclusion, the restrictions of the first six harmonics are
\[
\begin{array}{c|rrrrrr}
&H_0&H_1&H_2&H_3&H_4&H_5\\\hline
\mathbf1&1&0&0&0&0&0\\
D_+&0&1&0&0&0&1\\
D_-&0&0&0&1&0&1\\
A&0&0&0&1&1&0\\
B&0&0&1&0&1&1
\end{array}
\]
Interchanging the labels of the two five-cycle classes exchanges $D_+,D_-$ and changes no conclusion. The table follows from the torus weights $-j,-j+1,\ldots,j$, evaluated on rotations of orders two, three and five. In particular $H_2|_{A_5}=B$ and $H_5|_{A_5}=B\oplus C$.

The genuine rotation-group representations of dimensions six and twelve restricting to $P,Q$ are unique up to isomorphism. Indeed, their summands have dimensions $2j+1$, so only $j\leq2$ or $j\leq5$ can occur. Each module has a one-dimensional fixed subspace, so its rotation-group model contains exactly one copy of $H_0$. For $P$, the absence of $D_+$ forces the remaining five dimensions to be $H_2$. For $Q$, the absence of $A$ excludes $H_3,H_4$, and its $D_-$ constituent forces $H_5$. Thus the models used in Section F.3 are $H_0\oplus H_2$ and $H_0\oplus H_5$. This uniqueness concerns representation isomorphism classes, not a preferred identification with the geometric function spaces.

The matrix above has rank five, and its integral kernel is generated by the primitive vector
\[
(0,1,2,1,-1,-1),
\qquad
H_1+2H_2+H_3-H_4-H_5=0\quad\hbox{in }R(A_5).
\]
The last coordinate is $-1$, so a rational multiple of this vector with integral coordinates is already an integer multiple. Two reductions into $H_0,\ldots,H_5$ with the same restricted character consequently differ by just one integer parameter. Applying the fixed rule of Section F.3 independently to both columns of its harmonic table gives the same six multiplicities on both sides, so the difference between the resulting vectors is zero. This agreement is an additional calculation, not a consequence of equality in $R(A_5)$ alone. For instance $H_{11}$ and $2H_1\oplus2H_2\oplus H_3$ have identical restrictions but reduce to $(0,1,0,0,1,1)$ and $(0,2,2,1,0,0)$, respectively.

The common representation decomposes into the five irreducibles with multiplicities
\[
(46,156,156,199,245),
\]
and hence its equivariant endomorphism space has dimension
\[
46^2+2\cdot156^2+199^2+245^2=150414.
\]
The choices specified in Section F.3 give an isomorphism in \(\operatorname{Hom}_{A_5}(\Lambda^2\operatorname{Sym}^2Q,\Lambda^5\Lambda^2P)\).

\section{Level-five modular forms and character relations}\label{sec:modular}

\subsection{Cusps and weight-two modular forms}\label{subsec:G1}

Let $\Gamma(5)=\{\gamma\in\mathrm{SL}_2(\mathbb Z):\gamma\equiv I\pmod5\}$. A \emph{cusp} is a $\Gamma(5)$-orbit in $\mathbb Q\cup\{\infty\}$. Adjoining these boundary points to $\Gamma(5)\backslash\mathfrak H$ compactifies the quotient to $X(5)$. The image of $\Gamma(5)$ in $\mathrm{PSL}_2(\mathbb Z)$ has index sixty. Its cusp set is modeled by nonzero vectors of \(\mathbb F_5^2\) modulo sign, with twelve elements and a cyclic stabilizer of order five. Passing to their six underlying lines yields a two-to-one equivariant quotient. These give the same transitive permutation modules \(Q\) and \(P\) used earlier. The principal congruence group is torsion-free after passing to the projective action, and its genus is

\[
1+\frac{60}{12}-\frac{12}{2}=0.
\]

Duke \citep[Sections 3 and 6]{Duke2005} identifies the Rogers--Ramanujan function as a projective coordinate carrying the cusps to the icosahedral vertices. The manuscript uses the normalization of the degree-twelve divisor fixed in Section F.

A weight-two modular form is holomorphic on $\mathfrak H$, satisfies $f(\gamma\tau)=(c\tau+d)^2f(\tau)$ for $\gamma\in\Gamma(5)$, and is holomorphic at every cusp. Explicitly, if $\sigma=\left(\begin{smallmatrix}a&b\\c&d\end{smallmatrix}\right)\in\mathrm{SL}_2(\mathbb Z)$ sends infinity to that cusp, then $(c\tau+d)^{-2}f(\sigma\tau)$ has an expansion with no negative powers in the local parameter below. Since the projective action is torsion-free, the invariant differential $f(\tau)d\tau$ descends without interior singularities. Locally, after moving a cusp to infinity,

\[
u=e^{2\pi i\tau/5},\qquad d\tau=\frac5{2\pi i}\frac{du}{u}.
\]

Approaching the cusp means $u\to0$. Writing $f=a_0+a_1u+\cdots$, the differential has at most a simple pole, with residue $5a_0/(2\pi i)$. Conversely, every differential with at most simple poles at the cusps pulls back to such a form. Thus

\[
M_2(\Gamma(5))=H^0(X(5),K_{X(5)}(D)).
\]

Here $K_{X(5)}(D)$ allows meromorphic differentials with at most simple poles at $D$. The residue theorem and partial fractions identify this space equivariantly with the sum-zero subspace of \(\mathbb C[D]\); the kernel consists of holomorphic differentials and is zero because the genus is zero. Use the usual left action on differentials, \((g^{-1})^*\), so the residues transform by the induced permutation action. Consequently

\[
M_2(\Gamma(5))\cong B\oplus C,\qquad
\dim M_2(\Gamma(5))=11,\qquad S_2(\Gamma(5))=0.
\]

Cusp forms, denoted by $S_2$, have $a_0=0$ at every cusp, so their differentials are holomorphic and vanish on this genus-zero curve. Consequently \(S\cong M_2\oplus B\), \(\mathbf1\oplus S\cong M_2\oplus P\), and \(Q\cong\mathbf1\oplus M_2\) as finite-group representations. These isomorphisms concern the underlying \(A_5\)-representations.

\Needspace{12\baselineskip}
\section{The Mathieu subgroup and the doubled vertex representation}\label{sec:conformal}

\subsection{The marked subgroup and the doubled vertex space}\label{subsec:H1}

Use the standard degree-24 ATLAS permutations acting on \(\{1,\ldots,24\}\),

\[
\begin{aligned}
a={}&(1\,4)(2\,7)(3\,17)(5\,13)(6\,9)(8\,15)\\
&\quad(10\,19)(11\,18)(12\,21)(14\,16)(20\,24)(22\,23),\\
b={}&(1\,4\,6)(2\,21\,14)(3\,9\,15)(5\,18\,10)
(13\,17\,16)(19\,24\,23).
\end{aligned}
\]

In this subsection use right actions: \(x^{ab}=(x^a)^b\), so \(a\) acts first in the product \(ab\). Put

\begin{verbatim}
y = ababbababbababbababb
w = bbbbabbbbabbabbbbababbbaabbbbabababababb
z = w^(-1) y w.
\end{verbatim}

The ATLAS generators give the inclusion in \(M_{24}\) \citep{AtlasM24}. The subgroup \(\langle a,z\rangle\) has sixty elements, conjugacy-class sizes \(1,12,12,15,20\), and element-order counts \(1,15,20,24\) at orders \(1,2,3,5\). No proper union of these classes containing the identity has cardinality a nontrivial divisor of sixty, so the subgroup is simple and hence is isomorphic to \(A_5\).

Its action on \(\{1,\ldots,24\}\) has two orbits of length twelve. On each orbit the character values at the classes of orders \(1,2,3,5,5\) are \((12,0,0,2,2)\). Each stabilizer has order five, and the subgroups of order five in \(A_5\) are conjugate, so each orbit affords the vertex representation \(Q\). Writing \(\mathfrak a\) for the complex permutation representation on the twenty-four letters, we obtain
\[
\mathfrak a|_{A_5}\cong Q\oplus Q.
\]
An equivariant identification between the two orbits is given in the repository file \path{basis_data.json}. The cycle shapes on the twenty-four letters are \(1^{24},2^{12},3^8,1^4 5^4\), and the real fixed-point subspace is the two-dimensional space spanned by the two orbit sums.

\subsection{The doubled twelve-point module}

Using the restriction just proved, write \(\mathfrak a|_{A_5}\cong Q\otimes Z\), where \(Z=\mathbb C^2\) carries the trivial action. The Cauchy decomposition gives
\[
\Lambda^2(Q\otimes Z)
\cong
\bigl(\Lambda^2Q\otimes\operatorname{Sym}^2Z\bigr)
\oplus
\bigl(\operatorname{Sym}^2Q\otimes\Lambda^2Z\bigr)
\cong 3\Lambda^2Q\oplus\operatorname{Sym}^2Q.
\]
The dimensions are \(276=3(66)+78\). Thus the \(78\)-dimensional representation \(\operatorname{Sym}^2Q\) used in the \(3003\) identity is a summand of \(\Lambda^2\mathfrak a|_{A_5}\), the restriction of the weight-one space in the Conway construction of Duncan--Mack-Crane \citep{DuncanMackCrane2015}.

\section{Computational supplement and figures}\label{sec:reproducibility}

\subsection{Reproducing the finite calculations}\label{subsec:I1}

The repository contains exact-arithmetic implementations of the finite calculations used above. The principal programs are:

\begin{center}
\begin{tabular}{ll}
\toprule
calculation & program\\
\midrule
\(A_5\) characters and realization counts & \texttt{verify\_revision.py}\\
finite characters and signed orbits & \texttt{verify\_finite\_core.py}\\
faces and uniqueness at \(1540\) & \texttt{verify\_faces\_1540.py}\\
vertex triples at \(24310\) & \texttt{verify\_vertex\_triples.py}\\
cubic lines at \(7140\) & \texttt{verify\_character\_lines.py}\\
complete induction catalogue & \texttt{verify\_wave\_i.py}\\
manuscript and certificate agreement & \texttt{verify\_integration.py}\\
Klein and geometric calculations & \texttt{verify\_geometry.py}\\
Spin restrictions and arithmetic & \texttt{verify\_spin\_arithmetic.py}\\
marked \(M_{24}\) subgroup & \texttt{verify\_marked\_basis.py}\\
\bottomrule
\end{tabular}
\end{center}

Running
\[
\texttt{python run\_all.py}
\]
recomputes the principal finite calculations and rebuilds the manuscripts. The exact enumeration output is written to \path{verification/revision_results.json}, and the character, orbit, semisimple-moment, and residue data are stored under \path{verification/certificates/}. The analytic and classical results used in the article are cited at their point of use.

The five-coordinate order in the programs is $(\mathbf1,D_+,D_-,A,B)$. The file \path{verification/induction/wave_i_results.json} lists all $102$ nontrivial ordered pairs, their common characters, and their membership in the permutation, signed permutation, and complex monomial classes. The enumerators group representations by exterior character to find matching pairs.

The three construction programs and their JSON witnesses are in
\path{verification/constructions/}; the complete induction catalogue and
\texttt{wave\_i\_results.json} are in \path{verification/induction/}.
The integration check compares both manuscript tables with the exact
characters and verifies that the two independent enumerators return the
same complete lists of pairs. Use \texttt{python run\_all.py --all-math}
to include the specialist spinorial, arithmetic, marked-basis, modular,
and spectral calculations.

The programs in \path{verification/reconstruction/} independently recover
the subgroup characters, realization pairs, restricted counts, monomial actions of the representations constructed above, Klein weight reductions, and Mathieu word paths. They
are included in the standard command; only the final comparison program
reads the earlier certificates.

\subsection{The four figures}\label{subsec:I2}

Figure 1(a) shows the normalized half-triangle $2\leq k\leq n/2$, the omitted boundary and $k=1$ diagonal, and reflection $k\leftrightarrow n-k$. Figure 1(b) shows the local $3003$ pattern. Its entries and arrows are determined by the corresponding binomial coefficients and Pascal additions.

Figures 2 and 4 use the same labeled twelve-point orbit. Put
\[
\zeta=e^{2\pi i/5},\qquad \varphi=(1+\sqrt5)/2,
\]
and
\[
v_0=\infty,\quad v'_0=0,\quad
v_i=\varphi^{-1}\zeta^{i-1},\quad
v'_i=-\varphi\zeta^{i-1}\quad(1\leq i\leq5).
\]
These are the zeros of \(F\), paired by the antipodal map \(z\mapsto-1/\overline z\). Under inverse stereographic projection,
\[
z\longmapsto
\frac{(2\operatorname{Re}z,\,2\operatorname{Im}z,\,|z|^2-1)}
     {1+|z|^2},
\]
they form the vertices of a regular icosahedron: two poles and two pentagonal rings at heights \(\pm1/\sqrt5\). The six opposite pairs determine the axes used in Figure 2. The graph is the complete graph on the six axis labels, whose fifteen edges are their unordered pairs.

Figure 3 summarizes the compatible subgroup placement proved in Section D.4 and distinguishes the faithful spin lift from its projection to the orthogonal group. Figure 4 uses a chosen projective coordinate to identify the same twelve labels with the cusps of \(X(5)\); its circular arrangement is schematic, while the residue identification is proved in Section G.1.

\Needspace{8\baselineskip}
\begin{thebibliography}{19}
\setlength{\itemsep}{0.45em}
\bibitem{AtlasM24}
ATLAS of Finite Group Representations. \href{https://brauer.maths.qmul.ac.uk/Atlas/v3/permrep/M24G1-p24B0}{Degree-24 permutation representation of $M_{24}$}, generators M24G1-p24B0.g1 and M24G1-p24B0.g2. The exact permutations used here are transcribed in Section H.1.

\bibitem{Avanesov1967}
E. T. Avanesov. \href{https://eudml.org/doc/204811}{Solution of a problem on figurate numbers} (Russian). Acta Arithmetica 12(4) (1966/67), 409--420. DOI: 10.4064/aa-12-4-409-420.

\bibitem{Baez2026Coincidences}
John C. Baez. \href{https://johncarlosbaez.wordpress.com/2026/09/20/binomial-coefficient-coincidences/}{Binomial Coefficient Coincidences}. Azimuth, 20 September 2026.

\bibitem{BlokhuisBrouwerdeWeger2017}
Aart Blokhuis, Andries Brouwer, and Benne de Weger. \href{https://math.colgate.edu/~integers/r64/r64.pdf}{Binomial collisions and near collisions}. Integers 17 (2017), A64.

\bibitem{deWeger1997}
Benjamin M. M. de Weger. \href{https://math.deweger.net/papers/%5B20%5DdW-EqBinom-JNumTh%5B1997%5D.pdf}{Equal Binomial Coefficients: Some Elementary Considerations}. Journal of Number Theory 63(2) (1997), 373--386.

\bibitem{Duke2005}
W. Duke. \href{https://doi.org/10.1090/S0273-0979-05-01047-5}{Continued Fractions and Modular Functions}. Bulletin of the American Mathematical Society 42 (2005), 137--162.

\bibitem{DuncanMackCrane2015}
John F. R. Duncan and Sander Mack-Crane. \href{https://doi.org/10.1017/fms.2015.7}{The Moonshine Module for Conway's Group}. Forum of Mathematics, Sigma 3 (2015), e10. DOI: 10.1017/fms.2015.7.

\bibitem{Gallier2014Spin}
Jean Gallier. \href{https://arxiv.org/abs/0805.0311}{Clifford Algebras, Clifford Groups, and a Generalization of the Quaternions: The Pin and Spin Groups}. arXiv:0805.0311v3 (2014).

\bibitem{Lind1968}
D. A. Lind. \href{https://www.fq.math.ca/Scanned/6-3/lind.pdf}{The Quadratic Field $\mathbb Q(\sqrt5)$ and a Certain Diophantine Equation}. The Fibonacci Quarterly 6(3) (June 1968), 86--93.

\bibitem{LombardoSanders2010}
Sara Lombardo and Jan A. Sanders. \href{https://arxiv.org/abs/0912.1697}{On the Classification of Automorphic Lie Algebras}. Communications in Mathematical Physics 299 (2010), 793--824.

\end{thebibliography}
\end{document}
